% ======================================================================
% Paper 4 --- CONSTRAINTS: Admissibility Constraints and the Standard
% Model Field Content
%
% v1.0 (Feb 2026): PDF-only release (Zenodo DOI 10.5281/zenodo.18604845).
%   "Admissibility Constraints and Structural Saturation." Nine pages.
%   Focus: structural-constraint consequences of finite enforceability --
%   monogamy (T_M), area-law scaling, saturation (T4F: N_gen = 3),
%   horizons, overlap bounds (T25a/b: sin^2 theta_W constraint), dark
%   sector as capacity architecture, structural exclusions.
%
% v2.0 (April 2026): First .tex release. Full-scope rebuild absorbing
%   (a) all v1.0 content preserved and re-framed under PLEC; (b) the full
%   Standard Model field-content derivation previously PDF-only across
%   the codebase: gauge template uniqueness (Theorem_R), 1-of-4680 fermion
%   template selection, capacity partition C_total = 61 = 45 + 4 + 12,
%   generation count N_gen = 3 (T4F), weak mixing angle sin^2 theta_W
%   = 3/13 (T23/T24/T25/T27), mass structure (Schur chain), CKM + PMNS.
%
%   Five parts:
%     Part I --- Structural Constraints (preserves v1.0 content).
%     Part II --- Gauge Structure + Field Content + Generation Count.
%     Part III --- Mass, Mixing, and the Weak Mixing Angle.
%     Part IV --- Dark Sector, Relation to Quantum, Open Problems.
%     Part V --- Conclusion + forward references to Papers 5-7.
%
%   PLEC framing throughout. Epistemic tags [P] / [P_structural] /
%   [P]+[I] / [Pred] applied uniformly. Every theorem anchored to a
%   named check function via \coderef{name}{module.py}.
%
%   Accompanying Technical Supplement v1.0 contains full formal proofs
%   and red-team analysis.
% ======================================================================

\documentclass[11pt]{article}

\usepackage[T1]{fontenc}
\usepackage[utf8]{inputenc}
\usepackage[a4paper,margin=1in]{geometry}
\usepackage{setspace}
\usepackage{parskip}
\usepackage{enumitem}
\usepackage{booktabs}
\usepackage{array}
\usepackage{amsmath,amssymb,amsfonts}
\usepackage{amsthm}
\usepackage{mathtools}
\usepackage{graphicx}
\usepackage[most]{tcolorbox}
\usepackage{titlesec}
\usepackage{fancyhdr}

\titleformat{\part}[display]{\centering\Huge\bfseries}{\partname~\thepart}{1em}{}
\titleformat{\section}{\Large\bfseries}{\thesection.}{0.5em}{}
\titleformat{\subsection}{\large\bfseries}{\thesubsection}{0.5em}{}
\titleformat{\subsubsection}{\normalsize\bfseries}{\thesubsubsection}{0.5em}{}

\pagestyle{fancy}
\fancyhf{}
\fancyhead[L]{\small\itshape Paper 4: Constraints (v2.0)}
\fancyhead[R]{\small\itshape E.\,Brooke}
\fancyfoot[C]{\thepage}
\renewcommand{\headrulewidth}{0.4pt}

\usepackage[colorlinks=true,linkcolor=black,citecolor=black,urlcolor=black]{hyperref}

\tcbuselibrary{skins,breakable}

\tcbset{
    apfbox/.style={
        enhanced,
        colback=black!3,
        frame hidden,
        borderline={0.5pt}{0pt}{black},
        arc=0pt,
        left=8pt, right=8pt, top=8pt, bottom=8pt,
        fonttitle=\bfseries\color{black},
        coltitle=black,
        detach title,
        before upper={\tcbtitle\par\smallskip},
        before skip=14pt,
        after skip=14pt,
        grow to left by=-8pt,
        grow to right by=-8pt,
        sharp corners,
        breakable,
    }
}

\tcbset{
    wayfinder/.style={
        enhanced,
        colback=black!2,
        frame hidden,
        borderline={0.4pt}{0pt}{black!60},
        arc=0pt,
        left=10pt, right=10pt, top=10pt, bottom=10pt,
        fonttitle=\bfseries\color{black!80},
        coltitle=black!80,
        detach title,
        before upper={\tcbtitle\par\smallskip},
        before skip=16pt,
        after skip=16pt,
        sharp corners,
        breakable,
    }
}

\newcommand{\coderef}[2]{\smallskip\noindent\textit{Code reference:} \texttt{#1} in \texttt{#2}.}
\newcommand{\SU}{\mathrm{SU}}
\newcommand{\Un}{\mathrm{U}}
\newcommand{\epitag}[1]{\textbf{[#1]}}
\newcommand{\PLEC}{\emph{Principle of Least Enforcement Cost}}
\newcommand{\Aone}{\textbf{A1}}
\newcommand{\Atwo}{\textbf{A2}}
\newcommand{\MD}{\textbf{MD}}
\newcommand{\BW}{\textbf{BW}}
\newcommand{\paperref}[2]{\textbf{Paper~#1} (\textit{#2})}
\DeclareMathOperator*{\argmin}{arg\,min}

\newtheorem{theorem}{Theorem}[section]
\newtheorem{lemma}[theorem]{Lemma}
\newtheorem{proposition}[theorem]{Proposition}
\newtheorem{corollary}[theorem]{Corollary}
\theoremstyle{definition}
\newtheorem{definition}[theorem]{Definition}
\newtheorem{remark}[theorem]{Remark}

\title{
  \textbf{Paper 4: Constraints}\\
  \large Admissibility Constraints and the Standard Model Field Content\\
  \normalsize v2.0 --- April 2026\\
  \normalsize E.\,S.\,Brooke\\
  Admissibility Physics (APF) --- Canonical Series
}
\author{}
\date{}

\setstretch{1.05}
\setlist[itemize]{itemsep=3pt,topsep=4pt}
\setlist[enumerate]{itemsep=3pt,topsep=4pt}

\begin{document}
\maketitle

\begin{abstract}
\noindent
This paper derives the Standard Model field content --- the gauge group $\SU(3) \times \SU(2) \times \Un(1)$ with $N_c = 3$, the 45-fermion spectrum in three generations, and the associated mass and mixing structure --- from finite admissibility alone. The derivation rests on two foundations: (i) the structural consequences of finite enforceability established in \paperref{1}{Enforceability of Distinction} and \paperref{2}{Structure of Admissible Physics} --- monogamy of correlations, area-law scaling, saturation, and horizons --- and (ii) the capacity-partition architecture $C_{\mathrm{total}} = 61 = 45 + 4 + 12$ that follows from the admissible gauge template and the anomaly-free fermion filter.

The main structural constraints are derived first (Part~I). The capacity-saturation theorem $T_{4F}$ forces exactly three fermion generations: with $C_{\mathrm{int}} = \dim\SU(3) + \dim\SU(2) + \dim\Un(1) - d_{\mathrm{spacetime}} = 12 - 4 = 8$ and generation cost $E(n) = 2n + (n \bmod 2)$, the inequality $E(3) = 7 \leq 8$ is admissible while $E(4) = 9 > 8$ is not. The weak mixing angle is fixed at $\sin^2\theta_W = 3/13 \approx 0.2308$ (experimental: 0.2312, error 0.19\%) as the fixed point of the two-sector overlap/$\gamma$-ratio formula $T_{23}$. Dark matter is identified not as a particle species but as $C_{\mathrm{ext}}$ geometric correlation, with the cosmological constant $\Lambda$ emerging as the global capacity residual.

The Standard Model field content is then derived (Part~II). The gauge template is uniquely admissible (Theorem~R, anchored by $L_{\text{gauge\_template\_uniqueness}}$); the fermion spectrum is the unique survivor of 4,680 anomaly-and-Witten-compatible candidates; the capacity partition $61 = 45 + 4 + 12$ follows from the count of maintained distinctions. Mass structure, CKM, and PMNS mixing follow from capacity-per-dimension constraints and the generation-channel bridge (Part~III). The framework's predictions, open problems ($m_c$ at 2.6\%, absolute $m_t, m_b$ scales, neutrino mass hierarchy), and forward references to \paperref{5}{Quantum Structure}, \paperref{6}{Dynamics and Geometry}, and \paperref{7}{Action} close the paper (Part~IV).

Throughout, PLEC's four structurally necessary components --- \Aone{} (finite capacity), \MD{} (cost floor), \Atwo{} (argmin selection), \BW{} (non-degeneracy) --- appear as independent inputs, each doing specific work. Every claim traces to a named check function in the APF codebase v6.9, runnable via \texttt{python run\_checks.py}. Zero free parameters.
\end{abstract}

\vspace{0.5em}

\begin{tcolorbox}[wayfinder, title={Wayfinder: how to read this paper}]
\textbf{If you are looking for the Standard Model derivation:} start at Part~II (\S\ref{sec:gauge_template}). The three-generation derivation is \S\ref{sec:T4F}; the 1-of-4680 fermion filter is \S\ref{sec:anomaly_4680}; the weak mixing angle is \S\ref{sec:sin2thetaW}.

\textbf{If you are looking for the structural constraints (monogamy, area-law, saturation, horizons):} Part~I (\S\ref{sec:non_closure}--\S\ref{sec:horizons}). These are the foundation; Part~II builds on them.

\textbf{If you are looking for the dark-sector identification $\Lambda + \mathrm{DM} = C_{\mathrm{ext}}$:} \S\ref{sec:dark_sector} (Part~IV).

\textbf{If you are a reviewer:} the Technical Supplement (separate document) contains all formal proofs, the red-team section, and the full-detail countermodels. The main paper gives theorem statements + proof sketches $\leq$ 20 lines + \texttt{\textbackslash coderef} anchors.

\textbf{Epistemic conventions} (applied throughout): \epitag{P} = proved from PLEC; \epitag{P}+\epitag{I} = proved using an imported mathematical theorem (Frobenius, Brauer, Solèr, Gleason, HKM, Malament, Lovelock); \epitag{P\_structural} = proved modulo a sector-specific structural hypothesis named at the claim site; \epitag{Pred} = framework prediction awaiting experimental test. Every claim is anchored to a \texttt{check\_*} function in the v6.9 codebase via \texttt{\textbackslash coderef}.
\end{tcolorbox}

\tableofcontents
\vspace{1em}
\clearpage

% ======================================================================
\part{Structural Constraints}
% ======================================================================

% ======================================================================
\section{Introduction}\label{sec:intro}
% ======================================================================

\subsection{Constraints versus dynamics}\label{sec:intro_constraints_vs_dynamics}

\paperref{1}{Enforceability of Distinction}, \paperref{2}{Structure of Admissible Physics}, and \paperref{3}{Ledgers} established three load-bearing results of the framework: distinctions require enforcement against finite capacity (Paper~1); admissibility does not close on any single representation (Paper~2); dynamics and entropy track committed capacity (Paper~3). This paper asks the prior question that each of those assumes: \emph{what structural constraints must hold regardless of representation?}

Structural constraints are not dynamical laws that could be violated, statistical tendencies that could fluctuate, or approximations that could be refined. They are feasibility boundaries: configurations that exceed capacity, violate monogamy, overflow the admissible set, or cross saturation thresholds are inadmissible, full stop. Constraints apply even where dynamics, geometry, and probability cease to exist as representations. This is the key hierarchy: constraints are more fundamental than the formalisms built upon them.

\subsection{What this paper establishes}\label{sec:intro_what_paper_establishes}

Two classes of results, organized into five parts.

\textbf{Part~I --- Structural Constraints} (Sections~\ref{sec:inputs}--\ref{sec:horizons}). From PLEC alone, we derive: non-closure of admissibility under composition; monogamy of correlations ($T_M$, upgraded to \epitag{P} via biconditional proof); area-law scaling of entropy; saturation as a sharp feasibility boundary; horizons as regions of persistent saturation; overlap bounds constraining the weak mixing angle.

\textbf{Part~II --- Gauge Structure and Field Content} (Sections~\ref{sec:gauge_template}--\ref{sec:T4F}). The gauge template $\SU(3)\times\SU(2)\times\Un(1)$ with $N_c = 3$ is uniquely admissible (Theorem~R); the Standard Model fermion spectrum is the unique survivor of 4{,}680 anomaly-and-Witten-compatible candidates; the capacity partition $C_{\mathrm{total}} = 61 = 45 + 4 + 12$ follows structurally; the generation count $N_{\mathrm{gen}} = 3$ is forced by capacity saturation.

\textbf{Part~III --- Mass, Mixing, and the Weak Mixing Angle} (Sections~\ref{sec:sin2thetaW}--\ref{sec:pmns}). The weak mixing angle $\sin^2\theta_W = 3/13$ emerges as a two-sector fixed point; mass ratios follow a Schur structural chain; the Cabibbo angle and CKM matrix follow from the generation-channel bridge; PMNS lepton mixing and the seesaw structure are derived.

\textbf{Part~IV --- Dark Sector and Relation to Standard Physics} (Sections~\ref{sec:dark_sector}--\ref{sec:open_problems}). Dark matter is identified as $C_{\mathrm{ext}}$ geometric capacity; $\Lambda$ is the capacity residual; relation to standard quantum constraints (no-cloning, monogamy, Holevo, area laws) is made explicit; the open problems --- $m_c$ at 2.6\%, absolute $m_t/m_b$ scales, neutrino hierarchy --- are bounded.

\textbf{Part~V --- Closure} (Section~\ref{sec:conclusion}). Forward references to Papers~5--7 and the outstanding work items.

\subsection{Scope boundary}\label{sec:intro_scope_boundary}

\begin{tcolorbox}[apfbox, title={What Paper 4 does not claim}]
\begin{itemize}
    \item \emph{Absolute mass scales for $m_t$ and $m_b$} are Type-B anchors; absolute scale requires a $\sigma$ derivation that remains open.
    \item \emph{Neutrino mass hierarchy} is scoped via the Type-I seesaw (developed in \paperref{7}{Action}); specific hierarchy predictions remain open.
    \item \emph{Extra generations} ($N_{\mathrm{gen}} \geq 4$) are structurally excluded; hidden-generation speculation is not entertained.
    \item \emph{Dark-matter particle candidates} are not predicted; the framework commits to dark matter as geometric capacity ($C_{\mathrm{ext}}$), which is a falsifiable prediction (direct detection should find nothing).
    \item \emph{Dynamical dark energy} (DESI 2.8--4.2$\sigma$ preference for $w \neq -1$) is watched; framework commitment is $w = -1$ at structural level.
\end{itemize}
\end{tcolorbox}

% ======================================================================
\section{The admissibility constraint architecture}\label{sec:inputs}
% ======================================================================

\subsection{Inputs from Paper 1: PLEC's four components}\label{sec:inputs_plec}

This paper uses the four structurally necessary components of the \PLEC{} (PLEC), as established in \paperref{1}{Enforceability of Distinction} and its Technical Supplement:

\begin{itemize}
    \item \Aone{} (\emph{Finite Enforcement Capacity}): at any interface $\Gamma$, the total enforcement cost $\sum_{d \in G} \varepsilon(d) \leq C(\Gamma) < \infty$.
    \item \MD{} (\emph{Minimum Cost Floor}): non-trivial distinctions satisfy $\varepsilon(d) \geq \mu^* > 0$, uniformly across the distinction set.
    \item \Atwo{} (\emph{argmin selection}): the realized configuration is the argmin of total cost over the admissible set.
    \item \BW{} (\emph{Budget-window non-degeneracy}): the cost spectrum is rich enough that the argmin is structurally informative.
\end{itemize}

The four components are pairwise logically independent (Paper~1 Technical Supplement \S1, with explicit countermodels per component). The compact decomposition and each component's specific contribution is shown in Paper~0 Table~1 (\S3.5, \emph{PLEC component decomposition}), the \emph{PLEC component decomposition}, which is this paper's intended reference for ``what each component does.''

Two further PLEC-level pieces are used:

\begin{itemize}
    \item The \emph{Regime R} framework formalized in \texttt{apf/plec.py}: conditions R1 (smooth cost), R2 (local additivity), R3 (connected path space), R4 (no saturation) under which PLEC is well-posed as a variational problem.
    \item The five-type regime-exit taxonomy (Types I--V) formalizing where PLEC fails and what replaces it in each regime.
\end{itemize}
\coderef{check\_Regime\_R}{plec.py}

\subsection{Inputs from Paper 2: non-closure and the gauge template}\label{sec:inputs_paper2}

From \paperref{2}{Structure of Admissible Physics}:

\begin{itemize}
    \item \textbf{Non-closure theorem}: the admissibility algebra does not close on any single representation. No single Hilbert-space or $C^*$-algebra description covers all admissibility regimes.
    \item \textbf{Theorem R}: the gauge template $\SU(3) \times \SU(2) \times \Un(1)$ with $N_c = 3$ is uniquely admissible under PLEC. Three sub-steps (R1, R2, R3) establish $\SU(N_c)$ with $N_c \geq 3$ (non-abelian carrier), $\SU(2)$ (unique chiral carrier), and $\Un(1)$ with $N_c = 3$ (minimal completion).
    \item \textbf{Capacity-partition architecture}: the total capacity $C_{\mathrm{total}} = 61$ decomposes as $45 + 4 + 12$ (fermion distinctions + Higgs + gauge), determined by counting admissible distinctions in the gauge + field template.
\end{itemize}
\coderef{check\_Theorem\_R}{gauge.py}
\coderef{check\_L\_gauge\_template\_uniqueness}{gauge.py}
\coderef{check\_L\_count}{gauge.py}

These Paper~2 results are \emph{prerequisites} for Paper~4, not re-derived here. Paper~4 uses them as its own starting point and develops the structural-constraint and field-content consequences that follow.

\subsection{What Paper 4 derives}\label{sec:inputs_paper4_derives}

Everything in the following table is a Paper~4 result, anchored to a named check function in the v6.9 codebase:

\begin{center}\small
\begin{tabular}{@{}>{\raggedright\arraybackslash}p{0.30\linewidth} >{\raggedright\arraybackslash}p{0.42\linewidth} >{\raggedright\arraybackslash}p{0.20\linewidth}@{}}
\toprule
\textbf{Result} & \textbf{Content} & \textbf{Status} \\
\midrule
$T_M$ (Interface Monogamy) & Each distinction participates in at most one independent correlation & \epitag{P} \\
Area-law scaling & $S(A) \leq \kappa \cdot |\partial A|$; volume scaling inadmissible & \epitag{P} \\
Saturation (T4) & $S_\Gamma = C_\Gamma$ is a sharp feasibility boundary & \epitag{P} \\
Horizons (from persistent saturation) & Regions of permanent $S_\Gamma = C_\Gamma$; area-scaling forced & \epitag{P}+\epitag{I} \\
$N_{\mathrm{gen}} = 3$ (T4F) & $E(3) = 7 \leq 8 = C_{\mathrm{int}}$; $E(4) = 9 > 8$ & \epitag{P\_structural} \\
$\sin^2\theta_W = 3/13$ & Two-sector fixed point with overlap $x = 1/2$ and $\gamma_2/\gamma_1 = 17/4$ & \epitag{P\_structural} \\
1-of-4680 fermion filter & SM spectrum is the unique admissible survivor & \epitag{P} \\
45-fermion count (T\_field) & 15 per generation $\times$ 3 generations, anomaly-free & \epitag{P} \\
$C_{\mathrm{total}} = 61 = 45 + 4 + 12$ & Rigid capacity partition & \epitag{P} \\
Cabibbo $\sin\theta_C = e^{-3/2}$ & From generation-channel bridge & \epitag{P\_structural} \\
CKM structure (T\_CKM) & Element-by-element constraints from GCH & \epitag{P\_structural} \\
PMNS mixing (T\_PMNS) & Neutrino mixing from lepton-sector capacity & \epitag{P\_structural} \\
Dark matter = $C_{\mathrm{ext}}$ & Not a particle; geometric capacity residual & \epitag{P\_structural}+\epitag{Pred} \\
$\Lambda$ = capacity residual & $\Lambda \sim (C_{\mathrm{total}} - C_{\mathrm{used}}) / V$ & \epitag{P\_structural} \\
\bottomrule
\end{tabular}
\end{center}

The derivation order is: structural constraints first (Part~I), because they are regime-general and must be established before regime-specific field content; then the Standard Model content (Parts~II--III); then dark-sector identification and closure (Part~IV).

\subsection{Epistemic conventions}\label{sec:inputs_epistemic}

Every theorem statement is tagged:

\begin{itemize}
    \item \epitag{P} --- fully proved from PLEC and the prerequisites imported from Papers~1--3.
    \item \epitag{P}+\epitag{I} --- proved modulo an imported mathematical theorem (e.g., Lovelock uniqueness in $d=4$, Frobenius, Gleason, HKM). The external theorem is named at the claim site.
    \item \epitag{P\_structural} --- proved modulo a structural hypothesis specific to this paper or this sector, named at the claim site and flagged as a potential attack surface.
    \item \epitag{Pred} --- framework prediction awaiting experimental test. No framework-internal proof is claimed beyond structural consistency.
\end{itemize}

Proof depth in the main paper is at most a 20-line sketch per theorem, with the full formal proof in the Technical Supplement. Where a result has multiple layers, a compact tcolorbox states the theorem and the sketch follows. Every theorem has a \texttt{\textbackslash coderef} anchor to a named check function in \texttt{apf/*.py} (codebase v6.9).

\clearpage

% ======================================================================
\section{Non-closure under composition}\label{sec:non_closure}
% ======================================================================

\subsection{The fundamental property}\label{sec:non_closure_fundamental}

From \paperref{2}{Structure of Admissible Physics}: admissible correlation sets are not closed under composition. If $R_1$ and $R_2$ are individually admissible interfaces, $R_1 \cup R_2$ may be inadmissible. Joint enforcement cost may exceed joint capacity:
\[
E_\Gamma(R_1 \cup R_2) = E_\Gamma(R_1) + E_\Gamma(R_2) + I_\Gamma(R_1, R_2),
\]
where the interaction term $I_\Gamma$ captures the irreducible excess cost of joint enforcement. The excess is not an artifact of specific representations; it is a consequence of composition consuming capacity beyond the sum of individual pieces.

\coderef{check\_L\_nc}{core.py}

\subsection{Properties of non-closure}\label{sec:non_closure_properties}

Non-closure has several structural consequences, each independent and each load-bearing for Paper~4:

\textbf{Non-transitivity.} If $\Gamma_1$ enforces $A$--$B$ correlations and $\Gamma_2$ enforces $B$--$C$ correlations, $A$--$C$ correlations may nevertheless be inadmissible. Admissibility is interface-local and does not chain automatically.

\textbf{Path dependence.} The residual capacity after enforcing a sequence of distinctions depends on the order of enforcement, not merely the final set. Two sequences producing the same final configuration can consume different total capacity.

\textbf{No lattice structure.} Admissible sets do not form a lattice under union. This is a non-trivial obstruction: most resource theories in the quantum-information literature assume closure under free operations, which admissibility does not provide.

The absence of lattice structure is the formal expression of why \paperref{2}{Structure of Admissible Physics}'s non-closure theorem is structurally deep: no standard representation-theoretic approach can recover admissibility because no single lattice embedding exists.

\subsection{Non-closure as the foundation}\label{sec:non_closure_foundation}

Everything in Parts~I--III of this paper uses non-closure at some load-bearing step. Monogamy (\S\ref{sec:monogamy}) is non-closure applied to correlations sharing a common substrate. Saturation (\S\ref{sec:saturation}) is non-closure at the capacity boundary. The 1-of-4680 filter (\S\ref{sec:anomaly_4680}) uses non-closure to exclude fermion templates whose composition exceeds capacity. $N_{\mathrm{gen}} = 3$ (\S\ref{sec:T4F}) is non-closure of the $N_{\mathrm{gen}} = 4$ generation-count extension. Overlap bounds (\S\ref{sec:sin2thetaW}) are non-closure of mutually exclusive interface sectors.

Non-closure is the \emph{mechanism} that makes finite capacity physical, not a refinement of it. \Aone{} alone says the total is bounded; non-closure says the bound is sharp and that composition is where the bound bites.

\clearpage

% ======================================================================
\section{Monogamy of correlations}\label{sec:monogamy}
% ======================================================================

\subsection{The standard result, in quantum information}\label{sec:monogamy_standard}

For a tripartite quantum system $ABC$, entanglement between $A$ and the joint subsystem $BC$ bounds the sum of entanglements with $B$ and $C$ individually:
\[
E(A : BC) \;\geq\; E(A : B) + E(A : C).
\]
Various monogamy inequalities (CKW, squashed-entanglement variants, Holevo-type bounds) quantify this trade-off. In the standard presentation, monogamy is a feature of quantum mechanics, derived from Hilbert-space structure.

\subsection{The admissibility derivation}\label{sec:monogamy_admissibility}

In the admissibility framework, monogamy is not a quantum postulate. It is a direct consequence of finite interface capacity under non-closure.

Consider correlations enforced across interfaces. The interface $I_A$ between subsystem $A$ and the rest of the system (its complement, here $BC$) has finite capacity $C(I_A)$. Correlations with $B$ consume some capacity; correlations with $C$ consume more. When the total demand exceeds $C(I_A)$, joint enforcement fails. Monogamy is non-closure applied at a specific interface.

\subsection{Interface monogamy theorem}\label{sec:TM}

\begin{tcolorbox}[apfbox, title={Theorem $T_M$ (Interface Monogamy) \epitag{P}}]
Each distinction $d$ participates in at most one independent correlation. If $d$ participates in independent correlations with both $d_1$ and $d_2$, then $d_1$ and $d_2$ share the anchor $d$, making them not independent --- their enforcement budgets compete at $d$ (\Aone{}). Contradiction.
\end{tcolorbox}

\begin{proof}[Sketch]
Correlation requires both distinctions to exist as enforceable (definitional: one cannot correlate what is not there). Suppose $d$ participates in independent correlations with $d_1$ and $d_2$; independent means they do not share enforcement anchors. But both correlations are anchored at $d$. Therefore $d_1$ and $d_2$ do share the anchor $d$ --- they are not independent. Their enforcement budgets compete at $d$ (\Aone{}: finite capacity at $d$). Contradiction with independence. Therefore $d$ participates in at most one independent correlation. Full proof in the Technical Supplement, with the biconditional direction establishing that monogamy holds if and only if capacity is finite.
\end{proof}

\coderef{check\_T\_M}{supplements.py}

Theorem $T_M$, upgraded to \epitag{P} in the February~2026 release and maintained here, is the foundation for all subsequent monogamy-dependent results: overlap bounds on $\sin^2\theta_W$ ($T_{25a/b}$, \S\ref{sec:sin2thetaW}), saturation constraints, and the subordination bound $T_\eta \leq \varepsilon$. Previously treated as \epitag{P\_structural}, $T_M$ is now a biconditional proof from \Aone{} alone: monogamy holds if and only if capacity is finite.

\subsection{The capacity form of monogamy}\label{sec:monogamy_capacity_form}

\begin{tcolorbox}[apfbox, title={Theorem 3.1 (Monogamy inequality) \epitag{P}}]
For subsystems $A$, $B$, $C$ across context $\Gamma$:
\[
E_{\mathrm{sq}}(A : BC) \;\geq\; E_{\mathrm{sq}}(A : B) + E_{\mathrm{sq}}(A : C),
\]
where $E_{\mathrm{sq}}$ is the squashed correlation capacity --- the minimum irreducible load over all extensions. This is not a quantum postulate; it follows from finite capacity and non-closure.
\end{tcolorbox}

\subsection{Monogamy is capacity bookkeeping}\label{sec:monogamy_is_bookkeeping}

Monogamy is often presented as a peculiar feature of quantum mechanics, unique to the Hilbert-space formalism. Within the admissibility framework, it is simply capacity bookkeeping. The prediction: monogamy should appear in any system with interface-localized finite capacity, \emph{not only} quantum systems. The universality of monogamy across finite-capacity systems is a testable claim, distinct from any Hilbert-space axiomatization.

\clearpage

% ======================================================================
\section{Area-law scaling}\label{sec:area_law}
% ======================================================================

\subsection{The admissibility derivation}\label{sec:area_law_derivation}

Correlations between region $A$ and its complement must be enforced across the boundary $\partial A$. Enforcement capacity is localized at this interface. If capacity is finite per unit boundary, total enforceable correlation scales with boundary size:
\[
S(A) \;\leq\; \kappa \cdot |\partial A|.
\]
Area-law scaling is a direct consequence of interface-localized finite capacity. No quantum-gravitational argument is required; the area-law structure follows from \Aone{} applied to the boundary itself.

\coderef{check\_L\_area\_scaling}{core.py}

\subsection{Why not volume scaling}\label{sec:area_law_why_not_volume}

Volume scaling would require correlations to ``pass through'' the boundary repeatedly, each passage consuming capacity. Total demand would exceed interface capacity by a volume-to-area ratio that diverges as the region grows. Volume scaling is therefore inadmissible under \Aone{}.

This is the admissibility version of the holographic principle: entropy is a boundary phenomenon, not because of gravity, but because of interface-localized capacity. Gravity (derived in \paperref{6}{Dynamics and Geometry}) is one specific realization; the area-law structure precedes it.

\subsection{Regime dependence}\label{sec:area_law_regime}

The derivation assumes enforcement demand decomposes locally over extended interfaces --- a regime assumption (Regime R condition R2, \emph{local additivity}) valid when geometric representation is admissible. Near saturation, area scaling may break down. This breakdown reflects failure of local additivity (regime exit Type~II in the regime-exit taxonomy), not violation of admissibility itself. The framework predicts the regime where area scaling is exact and the corrections where it breaks down.

\clearpage

% ======================================================================
\section{Saturation as a feasibility boundary}\label{sec:saturation}
% ======================================================================

\subsection{Definition}\label{sec:saturation_definition}

\begin{tcolorbox}[apfbox, title={Definition 5.1 (Saturation)}]
Interface $\Gamma$ is \emph{saturated} when $S_\Gamma = C_\Gamma$. At saturation, no further independent correlations can be enforced across $\Gamma$.
\end{tcolorbox}

\subsection{Beyond saturation}\label{sec:saturation_beyond}

Saturation is a sharp feasibility boundary, not a smooth degradation. Below saturation, new correlations can be accommodated. At saturation, they cannot. Beyond saturation, the description framework itself fails: new structure requires different representations or reorganization of existing commitments. This is what regime-exit Type~IV captures: the boundary where PLEC's variational formulation ceases to be well-posed and a different regime takes over.

\coderef{check\_Regime\_exit\_Type\_IV}{plec.py}

\subsection{Capacity saturation forces three generations ($T_{4F}$)}\label{sec:T4F}

The most consequential application of saturation is the derivation of exactly three fermion generations.

\begin{tcolorbox}[apfbox, title={Theorem $T_{4F}$ (Capacity Saturation) \epitag{P\_structural}}]
The internal correlation capacity $C_{\mathrm{int}}$ derives from the gauge group dimensions:
\[
C_{\mathrm{int}} = \dim\SU(3) + \dim\SU(2) + \dim\Un(1) - d_{\mathrm{spacetime}} = 8 + 3 + 1 - 4 = 8.
\]
Each generation costs $E(n) = 2n + (n \bmod 2)$ units:
\begin{itemize}
    \item $N_{\mathrm{gen}} = 3$: $E(3) = 7 \leq 8$ \quad (admissible).
    \item $N_{\mathrm{gen}} = 4$: $E(4) = 9 > 8$ \quad (inadmissible).
\end{itemize}
Combined with $N_{\mathrm{gen}} \geq 3$ from CP violation (baryogenesis requires CKM complexity): $N_{\mathrm{gen}} = 3$ exactly.
\end{tcolorbox}

This is a saturation constraint in action: the fourth generation would exceed capacity. The proof is \emph{topological} --- it depends on routing structure, not numerical values. At $N_{\mathrm{gen}} \geq 4$, CP coherence necessarily fragments into multiple domains requiring partitioned enforcement capacity, which exceeds $C_{\mathrm{int}}$. The structural hypothesis tagged \epitag{P\_structural} is the generation-cost formula $E(n) = 2n + (n \bmod 2)$, justified in \paperref{2}{Structure} Supplement \S4 and re-derived in Paper~4's Technical Supplement \S5 from the generation-channel bridge.

\coderef{check\_T4F}{gauge.py}

\begin{center}\small
\begin{tabular}{@{}c c c l@{}}
\toprule
$N_{\mathrm{gen}}$ & $E(n)$ & $C_{\mathrm{int}} = 8$ & Status \\
\midrule
1 & 3 & $3 \leq 8$ & Admissible (but no CP violation) \\
2 & 5 & $5 \leq 8$ & Admissible (but no CP violation) \\
3 & 7 & $7 \leq 8$ & Admissible + CP violation \checkmark \\
4 & 9 & $9 > 8$ & Inadmissible \\
5 & 11 & $11 > 8$ & Inadmissible \\
\bottomrule
\end{tabular}
\end{center}

The argument is not ``the universe chose three generations'' but ``four or more generations violate $C_{\mathrm{int}} = 8$.'' The descriptive framing (Paper~0 \S5) applies: the realized generation count is located, not chosen.

\coderef{check\_T7}{gauge.py}

\subsection{Saturation is sharp, not gradual}\label{sec:saturation_sharp}

Saturation thresholds in the framework are discrete transitions. Below threshold: admissible. At threshold: marginally admissible (zero residual capacity). Above threshold: inadmissible, full stop. There is no soft degradation. The sharpness is a direct consequence of \Aone{}: a finite bound is sharp by definition; approach-to-saturation does not smoothly continue past the bound, because there is nothing beyond the bound within the framework's admissible set.

\clearpage

% ======================================================================
\section{Horizons from persistent saturation}\label{sec:horizons}
% ======================================================================

\subsection{Definition}\label{sec:horizon_definition}

\begin{tcolorbox}[apfbox, title={Definition 6.1 (Horizon)}]
Interface $\Gamma$ constitutes a \emph{horizon} when saturation is persistent: $S_\Gamma = C_\Gamma$ at all admissible continuations. No admissible path reduces the committed capacity.
\end{tcolorbox}

A horizon is saturation that does not relax. Individual transient saturations are ordinary feasibility boundaries; horizons are saturations embedded in the admissibility structure such that no continuation admitting them reduces $S_\Gamma$ below $C_\Gamma$.

\subsection{Consequences}\label{sec:horizon_consequences}

Across a horizon:
\begin{itemize}
    \item \textbf{New correlations are impossible.} Capacity is full; no further independent commitments can be made across the interface.
    \item \textbf{Signaling fails.} There is no capacity available for communication channels.
    \item \textbf{Geometric representation fails.} Smooth description exits its domain of validity. The horizon marks where smooth geometry ceases to be an admissible representation; a different regime (identifiable via the regime-exit taxonomy) takes over.
\end{itemize}

\subsection{Black hole horizons as the paradigmatic example}\label{sec:horizons_BH}

Black hole horizons are the paradigmatic physical example. Correlation capacity at the horizon scales with area (as expected from \S\ref{sec:area_law}):
\[
C_\Gamma \sim A/4\ell_P^2.
\]
Infalling matter commits capacity irreversibly. The Page transition --- the famous turnover in the black-hole entropy curve --- occurs when exterior correlations saturate horizon capacity; information is not lost, capacity is reorganized. The ``information paradox'' is resolved not by new quantum dynamics but by recognizing that the horizon is a saturation artifact and that capacity, once committed, cannot be both ``lost'' and ``preserved'' within admissible dynamics. Paper~6 develops the quantitative geometry; Paper~4 supplies the structural constraint.

\coderef{check\_T\_BH\_information}{supplements.py}
\coderef{check\_L\_BH\_page\_curve\_capacity}{supplements.py}

\clearpage

% ======================================================================
\section{Structural exclusions}\label{sec:exclusions}
% ======================================================================

The preceding sections derive the constraints. This section closes Part~I by listing what the constraints exclude --- configurations ruled out at the framework level, not by dynamical considerations.

\subsection{What cannot exist}\label{sec:exclusions_cannot_exist}

\begin{center}\small
\begin{tabular}{@{}>{\raggedright\arraybackslash}p{0.40\linewidth} >{\raggedright\arraybackslash}p{0.55\linewidth}@{}}
\toprule
\textbf{Excluded} & \textbf{Why} \\
\midrule
Unlimited entanglement & Would exceed interface capacity (\Aone{}) \\
Volume-law ground states & Boundary cannot support required correlations (\S\ref{sec:area_law}) \\
Perfect cloning & Would violate monogamy $T_M$ (\S\ref{sec:TM}) \\
Arbitrarily fine records & Saturation prevents unlimited refinement (\S\ref{sec:saturation}) \\
Global wavefunctions & Composition fails at horizons (\S\ref{sec:horizons}) \\
4th fermion generation & $E(4) = 9 > 8 = C_{\mathrm{int}}$ (\S\ref{sec:T4F}) \\
Dark-matter particles & Dark matter is $C_{\mathrm{ext}}$ geometric, not a species (\S\ref{sec:dark_sector}) \\
Arbitrarily many gauge groups & Non-closure caps gauge template at $\SU(3)\times\SU(2)\times\Un(1)$ \\
\bottomrule
\end{tabular}
\end{center}

Each exclusion is structural, not conventional. None is a choice the framework makes; each is forced by the interaction of \Aone{}, \MD{}, \Atwo{}, \BW{} with the specific substrate being described.

\subsection{What must exist}\label{sec:exclusions_must_exist}

\begin{center}\small
\begin{tabular}{@{}>{\raggedright\arraybackslash}p{0.40\linewidth} >{\raggedright\arraybackslash}p{0.55\linewidth}@{}}
\toprule
\textbf{Required} & \textbf{Why} \\
\midrule
Correlation trade-offs & Capacity is finite (\Aone{}) \\
Area-like scaling & Enforcement is interface-localized (\S\ref{sec:area_law}) \\
Saturation phenomena & Capacity exhaustion is possible (\S\ref{sec:saturation}) \\
Irreversibility & Commitment cannot be undone (Paper~3 $T_{\mathrm{second\_law}}$) \\
Exactly 3 generations & Forced by $T_{4F}$ (\S\ref{sec:T4F}) \\
SSB / Higgs mechanism & Unbroken vacuum inadmissible (anomaly cancellation) \\
Representational breakdown & Smooth descriptions have regime limits (\S\ref{sec:horizons}) \\
\bottomrule
\end{tabular}
\end{center}

Required features are not optional ingredients added to the theory. They are features whose absence would violate one of the framework's commitments. Removing irreversibility removes the second law; removing saturation phenomena removes finite capacity; removing the Higgs mechanism removes the possibility of broken vacua compatible with anomaly cancellation.

\clearpage

% ======================================================================
\part{The Standard Model Gauge Structure and Field Content}
% ======================================================================

% ======================================================================
\section{The gauge template: \texorpdfstring{$\SU(3)\times\SU(2)\times\Un(1)$}{SU(3) x SU(2) x U(1)} with \texorpdfstring{$N_c = 3$}{Nc = 3}}\label{sec:gauge_template}
% ======================================================================

Paper 2 established that the gauge template $\SU(3)\times\SU(2)\times\Un(1)$ with $N_c = 3$ is uniquely admissible. Paper~4 uses this template as input and develops its downstream consequences --- field content, generation count, weak mixing angle, mass and mixing structure. Because the gauge template is a prerequisite for everything in Parts~II and~III, we summarize the derivation chain before proceeding. The full formal proofs live in Paper~2's Technical Supplement; this section traces the chain and cites the codebase anchors.

\subsection{Theorem R and its three sub-steps}\label{sec:gauge_template_theoremR}

Theorem~R (\emph{Representation Requirements from Admissibility}, anchored by \texttt{check\_Theorem\_R} in \texttt{gauge.py}) proves that any admissible interaction theory satisfying \Aone{} must admit three distinct carrier representations, arrived at from carrier lemmas rather than assumed from a Lie-group classification:

\begin{itemize}
    \item \textbf{R1 (non-abelian, ternary carrier).} A faithful complex $N_c$-dimensional carrier admitting a trilinear invariant, with $N_c \geq 3$. The non-abelian requirement comes from $L_{\mathrm{nc}}$ (Paper~1); the trilinear invariant and $N_c \geq 3$ come from the stability of oriented composites under composition in a finite-capacity substrate.
    \item \textbf{R2 (pseudoreal chiral carrier).} A faithful pseudoreal 2-dimensional carrier. The pseudoreality is \emph{derived}: the 2-dimensional complex carrier on a finite-capacity interface admits a Witten mod-2 class obstruction unless equipped with a pseudoreal structure, which is what forces $\SU(2)$ rather than $\Un(2)$ or $\mathrm{SL}(2,\mathbb{C})$.
    \item \textbf{R3 (abelian grading).} A single abelian grading compatible with R1 and R2. Uniqueness of the abelian grading (only one $\Un(1)$ factor) follows from the enforcement-completeness argument rewritten in codebase v6.7: with four distinguishable representations supporting five physical states, exactly one $\Un(1)$ suffices; two $\Un(1)$ factors would provide an unused degree of freedom inadmissible under \Atwo{}'s cost-minimum selection.
\end{itemize}

\coderef{check\_Theorem\_R}{gauge.py}

\subsection{Gauge template uniqueness}\label{sec:gauge_template_uniqueness}

Theorem~R establishes that the three carrier requirements must be satisfied. $L_{\mathrm{gauge\_template\_uniqueness}}$ closes the gap between the abstract carrier requirements (Theorem~R) and the capacity-optimization argument (which selects $N_c = 3$ within the template) by proving that \emph{the template itself is forced}:

\begin{tcolorbox}[apfbox, title={$L_{\mathrm{gauge\_template\_uniqueness}}$ \epitag{P}}]
Given the three carrier requirements from Theorem~R, the gauge group must factor as
\[
G = \SU(N_c) \times \SU(2) \times \Un(1), \qquad N_c \geq 3 \text{ (odd)}.
\]
No alternative Lie-group structure satisfies all three requirements simultaneously.
\end{tcolorbox}

\begin{proof}[Sketch]
Six steps, all from \epitag{P} theorems: (i) R1 forces an $\SU(N_c)$ factor (simple, non-abelian, ternary carrier). (ii) R2 forces an $\SU(2)$ factor (the unique simple compact group with a faithful pseudoreal 2-dim representation). (iii) $N_c$ is odd: Witten-parity obstruction for even $N_c$. (iv) R3 forces a single $\Un(1)$ factor: multiple abelian gradings admit ambiguous hypercharge assignments inadmissible under \Atwo{}. (v) No further factor is admissible: additional simple factors increase $C_{\mathrm{int}}$ beyond what cost-minimization allows at the $N_c = 3$ solution. (vi) Direct-product structure (rather than a semi-direct product): the three factors commute at the substrate level because non-commutation would require a nested representation exceeding capacity (\Aone{}). Full proof in Paper~2 Technical Supplement \S4.
\end{proof}

\coderef{check\_L\_gauge\_template\_uniqueness}{gauge.py}

\subsection{Minimality and cost optimization: \texorpdfstring{$N_c = 3$}{Nc = 3}}\label{sec:gauge_template_minimality}

Theorem~R plus $L_{\mathrm{gauge\_template\_uniqueness}}$ narrows the template to $\SU(N_c)\times\SU(2)\times\Un(1)$ with $N_c \geq 3$ odd. The \emph{specific} value $N_c = 3$ is then determined by capacity-budget optimization:

\begin{tcolorbox}[apfbox, title={$T_{\mathrm{gauge}}$ (SU(3)$\times$SU(2)$\times$U(1) by cost optimization) \epitag{P}}]
At the \Atwo{} argmin of total enforcement cost over admissible gauge templates satisfying Theorem~R, the unique solution is
\[
(N_c, \text{chiral sector, grading}) = (3, \SU(2), \Un(1)),
\]
with internal capacity $C_{\mathrm{int}} = \dim \SU(3) + \dim \SU(2) + \dim \Un(1) - d_{\mathrm{spacetime}} = 8 + 3 + 1 - 4 = 8$.
\end{tcolorbox}

\begin{proof}[Sketch]
$N_c = 3$ emerges from the cubic anomaly equation solved per candidate $N_c$ (not hardcoded): for each odd $N_c \geq 3$, the cubic-anomaly-cancellation condition $[\SU(N_c)]^3 = 0$ combined with the minimal admissible fermion spectrum gives a total enforcement cost; $N_c = 3$ minimizes this cost. $N_c = 5, 7, \ldots$ produce strictly higher-cost admissible sectors that lose at the \Atwo{} argmin. The codebase solves this optimization explicitly: \texttt{check\_T\_gauge} runs the anomaly equation for each $N_c$ and reports the unique cost-minimum winner.
\end{proof}

\coderef{check\_T\_gauge}{gauge.py}

This is the foundation on which the 45-fermion spectrum of \S\ref{sec:anomaly_4680} and the capacity partition $C_{\mathrm{total}} = 61$ of \S\ref{sec:capacity_61} both build. The template $\SU(3)\times\SU(2)\times\Un(1)$ with $N_c = 3$ and $C_{\mathrm{int}} = 8$ is the fixed substrate for everything that follows.

% ======================================================================
\section{Anomaly cancellation and the 1-of-4680 fermion filter}\label{sec:anomaly_4680}
% ======================================================================

The gauge template of \S\ref{sec:gauge_template} is the substrate; it does not yet fix the fermion content. The Standard Model's specific fermion spectrum --- 15 chiral Weyl spinors per generation, carrying the charges $\{Q(3,2), L(1,2), u(\bar 3, 1), d(\bar 3, 1), e(1,1)\}$ --- is selected from a much larger candidate pool of gauge-representation assignments. This section derives that selection.

The logic runs in the \emph{opposite} direction from standard physics. In standard QFT, anomaly cancellation is imposed as a consistency requirement: any chiral gauge theory must be anomaly-free to preserve unitarity and renormalizability, and the fermion content is then \emph{chosen} to satisfy the cancellation conditions. In the admissibility framework, the logic is:

\begin{enumerate}
    \item The gauge group $\SU(3)\times\SU(2)\times\Un(1)$ with $N_c = 3$ is derived from capacity optimization ($T_{\mathrm{gauge}}$, \S\ref{sec:gauge_template_minimality}).
    \item The fermion template is scanned exhaustively over all admissible representation assignments.
    \item Anomaly cancellation, Witten parity, and the framework's capacity-based filters are imposed.
    \item Exactly one template --- the Standard Model's --- survives.
\end{enumerate}

The 1-of-4680 result is thus not an imposed consistency condition; it is a derivation outcome.

\subsection{Anomaly-freedom as an admissibility constraint}\label{sec:anomaly_free}

In the admissibility framework, gauge anomalies are not a quantum artifact to be avoided by consistency; they are the signature of capacity overflow in a proposed gauge-matter coupling. A gauge theory with uncancelled anomalies attempts to enforce a distinction structure that exceeds interface capacity at the anomaly node. The seven gauge anomaly conditions (triangle anomalies for $[\SU(3)]^3$, $[\SU(2)]^2 \Un(1)$, $[\Un(1)]^3$, mixed gravitational, and Witten mod-2) are each an \Aone{} constraint evaluated at a specific composite enforcement interface.

\begin{tcolorbox}[apfbox, title={$L_{\mathrm{anomaly\_free}}$ \epitag{P}}]
The framework-derived particle content with framework-derived hypercharges satisfies all seven gauge anomaly cancellation conditions per generation and for $N_{\mathrm{gen}} = 3$ generations combined. The cancellation is \emph{exact, not accidental}: each condition is an \Aone{} capacity equality at its specific interface.
\end{tcolorbox}

\coderef{check\_L\_anomaly\_free}{gauge.py}

This result is a structural cross-check: the derivation that produced the particle content and the derivation that imposes anomaly cancellation yield the same outcome. That agreement is not a coincidence but a consequence of both derivations flowing from \Aone{} at their respective interfaces.

\subsection{Witten parity}\label{sec:witten_parity}

The Witten global anomaly is a $\mathbb{Z}_2$ obstruction: an $\SU(2)$ gauge theory with an odd number of fundamental doublets is inadmissible because a large gauge transformation ($\pi_4(\SU(2)) = \mathbb{Z}_2$) changes the sign of the fermion determinant, producing a non-gauge-invariant path integral. In standard QFT this is an anomaly to be avoided; in the admissibility framework, it is a non-closure constraint: an odd number of $\SU(2)$ doublets produces a composite enforcement interface that cannot be consistently defined across the $\pi_4$ loop.

The Standard Model has an even number of $\SU(2)$ doublets per generation (two, from $Q$ and $L$), and the 3-generation total is even (six). Witten parity is satisfied.

\coderef{check\_L\_Witten\_parity}{gauge.py}

Witten parity combined with the seven triangle anomaly conditions gives eight independent constraints on the fermion assignment. These define the filter that the 4{,}680-candidate scan applies.

\subsection{The 4{,}680-candidate scan}\label{sec:4680_scan}

The exhaustive derivation of the SM fermion content proceeds in two phases, implemented directly in \texttt{check\_T\_field}:

\textbf{Phase 1: Exhaustive scan.} 4{,}680 fermion templates are constructed by taking all admissible combinations of:
\begin{itemize}
    \item $\SU(3)$ representations from $\{3, \bar 3, 6, \bar 6, 8\}$ (the admissible low-dimensional reps under $C_{\mathrm{int}} = 8$).
    \item $\SU(2)$ representations from $\{1, 2\}$ (singlet + doublet; higher reps excluded by \MD{}+\Aone{} at the internal-capacity budget).
    \item Up to 5 distinct field types.
    \item Up to 3 colored singlets and 2 lepton singlets (the capacity-compatible admissible bounds).
\end{itemize}

Seven filters are applied sequentially: anomaly-freedom for $\SU(3)$ and $\SU(2)$, chirality (left-handed Weyl), the $[\SU(3)]^3$ cubic anomaly condition, Witten parity, the full triangle-anomaly set, and the CPT quotient. Minimality under \Atwo{} then selects the unique cost-minimum winner.

\textbf{Phase 2: Closed-form exclusions.} Five categories outside the Phase~1 scan are ruled out by closed-form arguments:
\begin{itemize}
    \item $\SU(3)$ representations of dimension $\geq 10$ (i.e., $10, \bar{10}, 15, \ldots$): a single field in such a representation exceeds $C_{\mathrm{int}}$.
    \item Higher $\SU(2)$ representations (triplet, quartet, \ldots): fail both chirality and the anomaly-cancellation filter.
    \item More than 5 distinct field types per generation: violates \Atwo{} minimality given the capacity partition.
    \item Vector-like (non-chiral) matter: fails the chirality filter; vector-like pairs annihilate to gauge-invariant trivial content.
    \item Exotic gradings (multiple $\Un(1)$ factors): ruled out by $L_{\mathrm{gauge\_template\_uniqueness}}$'s uniqueness of the abelian grading.
\end{itemize}

\coderef{check\_T\_field}{gauge.py}

\subsection{The unique admissible survivor}\label{sec:4680_survivor}

Both phases converge on a single winner:

\begin{tcolorbox}[apfbox, title={$T_{\mathrm{field}}$ (SM fermion content) \epitag{P}}]
Given $\SU(3)\times\SU(2)\times\Un(1)$ with $N_c = 3$ ($T_{\mathrm{gauge}}$) and $N_{\mathrm{gen}} = 3$ ($T_7$), the Standard Model fermion template
\[
\{Q(3, 2)_{1/6}, \; L(1, 2)_{-1/2}, \; u(\bar 3, 1)_{-2/3}, \; d(\bar 3, 1)_{1/3}, \; e(1, 1)_{+1}\}
\]
is the unique chiral fermion content satisfying all admissibility constraints. $15$ Weyl spinors per generation $\times 3$ generations $= 45$ Weyl fermions total.
\end{tcolorbox}

The ``1-of-4{,}680'' is a statement about the filter ratio: of the 4{,}680 gauge-representation templates scanned in Phase~1, exactly one survives the seven filters under \Atwo{} minimality. That one is the SM. No other template passes. The Phase~2 closed-form exclusions show that no template \emph{outside} the 4{,}680-candidate pool could survive either --- so ``1-of-4{,}680'' understates the uniqueness: it is 1-of-all-admissible.

\coderef{check\_T\_field}{gauge.py}

% ======================================================================
\section{The 45-fermion spectrum and the capacity partition \texorpdfstring{$C_{\mathrm{total}} = 61$}{Ctotal = 61}}\label{sec:capacity_61}
% ======================================================================

The 1-of-4680 result (\S\ref{sec:anomaly_4680}) fixes the fermion content at 45 Weyl spinors. Combined with the gauge template of \S\ref{sec:gauge_template} and the Higgs scalar content, the total enforcement-capacity count $C_{\mathrm{total}} = 61$ emerges as a rigid partition. This section derives the partition and explains its cosmological consequences (full quantitative cosmological development is Paper~6's territory; Paper~4 establishes the partition).

\subsection{The 45-fermion count}\label{sec:45_fermion}

The 45 Weyl fermions decompose per generation as follows. Each generation has 15 distinct Weyl-spinor degrees of freedom, grouped into five multiplets:

\begin{center}\small
\begin{tabular}{@{}l l c c l@{}}
\toprule
\textbf{Multiplet} & \textbf{$(\SU(3), \SU(2))$} & \textbf{$Y$} & \textbf{Weyl DOF} & \textbf{Content} \\
\midrule
$Q$ & $(3, 2)$    & $+1/6$ & $6$ & Left-handed quark doublet (u, d) in three colors \\
$L$ & $(1, 2)$    & $-1/2$ & $2$ & Left-handed lepton doublet ($\nu$, $e$) \\
$u^c$ & $(\bar 3, 1)$ & $-2/3$ & $3$ & Right-handed up-type quark singlet \\
$d^c$ & $(\bar 3, 1)$ & $+1/3$ & $3$ & Right-handed down-type quark singlet \\
$e^c$ & $(1, 1)$    & $+1$ & $1$ & Right-handed charged lepton singlet \\
\midrule
\multicolumn{4}{r}{Per-generation total:} & 15 Weyl DOF \\
\multicolumn{4}{r}{3-generation total:} & \textbf{45 Weyl DOF} \\
\bottomrule
\end{tabular}
\end{center}

Right-handed neutrinos are not in the framework's chiral content; the seesaw mechanism (Paper~7) handles neutrino masses via Majorana mass terms on heavy right-handed singlets that are generated dynamically rather than present in the gauge-charged fermion content. This is consistent with the 1-of-4{,}680 scan, which admits only chiral, gauge-charged Weyl spinors; Majorana singlets contribute to the seesaw but not to the primary capacity count.

\coderef{check\_T\_field}{gauge.py}

\subsection{Capacity partition \texorpdfstring{$61 = 45 + 4 + 12$}{61 = 45 + 4 + 12}}\label{sec:61_partition}

The Standard Model's total structural-enforcement count partitions as:

\begin{tcolorbox}[apfbox, title={$L_{\mathrm{count}}$ (Capacity counting) \epitag{P}}]
At Bekenstein saturation, the number of independently enforceable capacity units equals the number of \emph{structural enforcement channels}: one per chiral species, one per gauge-automorphism direction, one per Higgs real component. Kinematic degrees of freedom (polarizations, helicities) do not contribute independent channels. The partition is:
\[
C_{\mathrm{total}} \;=\; 45 \; (\text{fermions}) \;+\; 4 \; (\text{Higgs}) \;+\; 12 \; (\text{gauge}) \;=\; 61.
\]
\end{tcolorbox}

The three summands:
\begin{itemize}
    \item \textbf{45 fermion channels:} from $T_{\mathrm{field}}$ (\S\ref{sec:45_fermion}).
    \item \textbf{4 Higgs channels:} the complex doublet $H \in (1, 2)_{1/2}$ carries 4 real scalar degrees of freedom. Three are absorbed by the Goldstone mechanism as the longitudinal components of $W^\pm, Z^0$ (contributing no independent channel after electroweak symmetry breaking but counting pre-SSB); one remains as the physical Higgs boson. At the \emph{capacity} level --- which is counted pre-SSB --- all four count. Post-SSB counting gives $4 = 3_{\mathrm{Goldstone}} + 1_{\mathrm{Higgs}}$; the distinction is dynamical, not structural.
    \item \textbf{12 gauge channels:} $\dim \SU(3) + \dim \SU(2) + \dim \Un(1) = 8 + 3 + 1 = 12$. These are the gauge-automorphism directions, each an independent capacity-commitment channel.
\end{itemize}

\coderef{check\_L\_count}{gauge.py}

The partition is rigid under the admissibility framework: removing any of the three summands violates anomaly cancellation, chirality, or gauge-group uniqueness. Changing the 45 requires changing the 1-of-4680 scan; changing the 4 requires a different Higgs content (and the framework excludes two-Higgs-doublet models at this scope); changing the 12 requires a different gauge group (excluded by Theorem~R + $L_{\mathrm{gauge\_template\_uniqueness}}$). The count is not adjustable.

\subsection{Partition structure and cosmological consequences}\label{sec:61_cosmological}

The capacity partition has direct cosmological consequences: with $C_{\mathrm{total}} = 61$ partitioned as above, the admissibility-level prediction for cosmological density parameters emerges from the count of capacity committed to each structural role. At Planck-scale saturation, the framework derives:

\begin{center}\small
\begin{tabular}{@{}l c c l@{}}
\toprule
\textbf{Parameter} & \textbf{APF prediction} & \textbf{Planck 2018} & \textbf{Structural role} \\
\midrule
$\Omega_b$         & $3/61 \approx 0.0492$  & $0.0493 \pm 0.0006$ & Color-singlet baryonic capacity \\
$\Omega_c$         & $16/61 \approx 0.2623$ & $0.2607 \pm 0.0062$ & Color-charged (confined) capacity \\
$\Omega_m$         & $19/61 \approx 0.3115$ & $0.3111 \pm 0.0056$ & Total matter capacity \\
$\Omega_\Lambda$   & $42/61 \approx 0.6885$ & $0.6889 \pm 0.0056$ & Geometric capacity residual \\
\bottomrule
\end{tabular}
\end{center}

All four match Planck 2018 within 1.2\%. The full cosmological derivation (including the four-parameter match and the $H_0 = 67.76$ km/s/Mpc prediction) is in Paper~6 \S11. Paper~4's contribution is the partition architecture: $45 + 16 + 3 + 12 \to 61$ with the matter/radiation/geometric groupings that Paper~6 quantifies.

\coderef{check\_L\_saturation\_partition}{cosmology.py}

% ======================================================================
\section{Generation count \texorpdfstring{$N_{\mathrm{gen}} = 3$}{Ngen = 3} by capacity ladder}\label{sec:Ngen3_ladder}
% ======================================================================

The generation count $N_{\mathrm{gen}} = 3$ has two derivations in the framework, reaching the same conclusion via different paths. The topological-saturation argument ($T_{4F}$, \S\ref{sec:T4F}) uses the cost formula $E(n) = 2n + (n \bmod 2)$ applied to the internal capacity $C_{\mathrm{int}} = 8$. A second, complementary route uses the capacity-ladder framework of \texttt{generations.py}, where each generation is a distinct ``rung'' on a charge ladder and ladder-termination follows from Froggatt--Nielsen uniqueness combined with \BW{} non-degeneracy. This section develops the capacity-ladder route and relates it to $T_{4F}$.

\subsection{The capacity-ladder framework}\label{sec:ladder_framework}

The capacity ladder (\texttt{check\_T\_capacity\_ladder}) assigns to each generation $g \in \{1, 2, 3, \ldots\}$ a capacity charge $q_g$ drawn from the admissible enforcement budget:

\begin{tcolorbox}[apfbox, title={$T_{\mathrm{capacity\_ladder}}$ \epitag{P}}]
Generations are ordered by increasing capacity charge: $q_1 < q_2 < q_3 < \ldots$. Each rung $q_g$ is a distinct admissible value of the Froggatt--Nielsen (FN) ladder-spacing operator. The admissible ladder terminates when the cumulative capacity $\sum_{g=1}^{N_{\mathrm{gen}}} E(q_g)$ exceeds $C_{\mathrm{int}}$.
\end{tcolorbox}

\coderef{check\_T\_capacity\_ladder}{generations.py}

The ladder framework is a discrete version of the continuous capacity-budget calculation underlying $T_{4F}$. Where $T_{4F}$ uses a single aggregate cost formula $E(n) = 2n + (n \bmod 2)$, the ladder framework decomposes the cost per generation and sums: each generation contributes a specific increment, and the ladder terminates at the first rung whose cumulative sum exceeds $C_{\mathrm{int}} = 8$.

\subsection{FN ladder uniqueness}\label{sec:generation_cost_formula}

The Froggatt--Nielsen ladder structure is not a postulate; it is the unique admissible decomposition of the generation charge assignment:

\begin{tcolorbox}[apfbox, title={$L_{\mathrm{FN\_ladder\_uniqueness}}$ \epitag{P\_structural}}]
Given the capacity partition $C_{\mathrm{total}} = 61$ and the gauge template of \S\ref{sec:gauge_template}, the admissible generation-charge assignment is a strictly increasing sequence $q_1 < q_2 < q_3$ with rung spacings uniquely determined by the FN flavor-symmetry structure. Any alternative decomposition (equal charges, decreasing charges, non-monotonic) violates \BW{} non-degeneracy at the generation level.
\end{tcolorbox}

\coderef{check\_L\_FN\_ladder\_uniqueness}{generations.py}

The structural hypothesis (\epitag{P\_structural}) is the identification of the FN flavor symmetry structure with the capacity-charge ladder; this is a sector-specific assumption named here as a potential attack surface. The Technical Supplement \S6 develops the full FN ladder analysis and the justification for the identification.

\subsection{Comparison with the topological \texorpdfstring{$T_{4F}$}{T4F} argument}\label{sec:T4F_vs_ladder}

The topological argument (\S\ref{sec:T4F}) and the capacity-ladder argument (this section) are two routes to the same conclusion:

\begin{center}\small
\begin{tabular}{@{}l l l@{}}
\toprule
\textbf{Route} & \textbf{Mechanism} & \textbf{Where $N_{\mathrm{gen}} = 4$ fails} \\
\midrule
$T_{4F}$ (topological) & Aggregate $E(n) = 2n + (n \bmod 2)$ & $E(4) = 9 > 8 = C_{\mathrm{int}}$ \\
Capacity ladder & Incremental FN ladder sum & Cumulative sum at rung 4 exceeds $C_{\mathrm{int}}$ \\
\bottomrule
\end{tabular}
\end{center}

The two routes agree at $N_{\mathrm{gen}} = 3$ (both admissible) and $N_{\mathrm{gen}} = 4$ (both inadmissible). The ladder framework is the more granular derivation, because it explains \emph{why} each individual generation is admissible: the first rung $q_1$ is free (lowest admissible charge), the second rung $q_2$ fits in the remaining budget, the third rung $q_3$ exhausts the budget near-saturation, and the fourth rung overflows. The topological argument is the compact summary.

Both routes are \epitag{P\_structural}: both rely on the structural hypothesis that the per-generation cost follows the specific formula ($T_{4F}$) or the FN ladder decomposition (capacity ladder). The Technical Supplement \S5--6 develops both in detail.

This completes Part~II: the gauge template, the 45-fermion field content, the capacity partition $C_{\mathrm{total}} = 61$, and the generation count $N_{\mathrm{gen}} = 3$ are all derived. Part~III turns to mass, mixing, and the weak mixing angle.

\clearpage

% ======================================================================
\part{Mass, Mixing, and the Weak Mixing Angle}
% ======================================================================

% ======================================================================
\section{The weak mixing angle: \texorpdfstring{$\sin^2\theta_W = 3/13$}{sin2 thetaW = 3/13}}\label{sec:sin2thetaW}
% ======================================================================

The weak mixing angle $\theta_W$ is the rotation between the electroweak $\SU(2)_L \times \Un(1)_Y$ gauge eigenstates and the physical $W^\pm, Z^0, \gamma$ mass eigenstates. In the Standard Model, $\sin^2\theta_W$ is a free parameter fit to experiment: $\sin^2\theta_W \approx 0.23122(4)$ at the $Z$ pole. In the admissibility framework, $\sin^2\theta_W$ is derived structurally as the fixed point of a two-sector interface-capacity equilibrium, yielding the clean rational value $3/13 \approx 0.2308$, which agrees with experiment to 0.19\%.

The derivation proceeds in four steps. First, the interface monogamy theorem $T_M$ bounds the overlap parameter $x$ between the electroweak sectors ($T_{25a}$). Second, saturation tightens the bound to $x = 1/2$ ($T_{25b}$). Third, representation counting gives the gamma ratio $\gamma_2/\gamma_1 = 17/4$ ($T_{27d}$). Fourth, the two-sector fixed-point formula $T_{23}$ combines $x$ and $\gamma$ into a closed-form prediction. Gate $S_0$ (interface-schema invariance) ensures the result is basis-independent.

\subsection{Interface monogamy bounds ($T_{25a}$)}\label{sec:T25a}

The electroweak $\Un(1)_Y$ hypercharge and $\SU(2)_L$ weak isospin are enforced at interfaces sharing partial capacity. Interface monogamy ($T_M$, \S\ref{sec:TM}) bounds the overlap parameter $x$ between the $m = 3$ electroweak mixing channels ($W^+, W^-, Z/\gamma$) --- more generally, $m$ mixer channels plus 1 bookkeeper channel $= 4$ EW channels at $C_{\mathrm{EW}} = \kappa \cdot 4 = 8$, constraining $x$ by the cut-set argument:

\begin{tcolorbox}[apfbox, title={$T_{25a}$ (Interface monogamy bound on $x$) \epitag{P}}]
\[
x \;\in\; \left[\frac{1}{m}, \frac{m-1}{m}\right] \;=\; \left[\frac{1}{3}, \frac{2}{3}\right].
\]
Each sector must contribute at least $1/m$ overlap by the cut-set argument; symmetrically, no sector can exceed $(m-1)/m$.
\end{tcolorbox}

\coderef{check\_T25a}{generations.py}

\subsection{Saturation tightening ($T_{25b}$)}\label{sec:T25b}

Near-saturation ($T_{4F}$: $7/8 = 87.5\%$ capacity utilization at $N_{\mathrm{gen}} = 3$), the overlap parameter is driven toward its symmetric midpoint. If $x$ deviates far from $1/2$, one sector overflows while another underuses capacity; at saturation, the monogamy bound and capacity efficiency jointly force:

\begin{tcolorbox}[apfbox, title={$T_{25b}$ (Saturation tightening) \epitag{P\_structural}}]
At EW saturation,
\[
x \;=\; \frac{1}{2}.
\]
\end{tcolorbox}

\coderef{check\_T25b}{generations.py}

The structural hypothesis is that EW capacity is effectively saturated --- which follows from $T_{4F}$'s capacity utilization of $7/8$ but relies on the identification of EW saturation with generation-sector saturation. Full justification in Technical Supplement \S7.

\subsection{The gamma ratio from representation counting ($T_{27d}$)}\label{sec:T27d}

The second input to the fixed-point formula is the ratio $\gamma_2/\gamma_1$, where $\gamma_i$ are the gauge-capacity ``braiding'' coefficients for the $\SU(2)$ and $\Un(1)$ sectors respectively. Representation counting at the gauge-capacity fixed point gives:

\begin{tcolorbox}[apfbox, title={$T_{27d}$ (Gamma ratio) \epitag{P\_structural}}]
\[
\frac{\gamma_2}{\gamma_1} \;=\; \frac{17}{4}.
\]
\end{tcolorbox}

\coderef{check\_T27d}{generations.py}

The ratio $17/4$ emerges from counting admissible gauge-capacity configurations weighted by the multiplicity of representations at each sector --- the $\SU(2)$ side admits 4 admissible charge configurations for the mixer triple, and the $\Un(1)$ side admits 17 admissible gauge-capacity assignments at the fixed point. The full counting argument is in Technical Supplement \S7.3.

\subsection{The two-sector fixed-point formula ($T_{23}$)}\label{sec:T23}

The three inputs $(x, \gamma_1, \gamma_2)$ combine through the two-sector fixed-point formula:

\begin{tcolorbox}[apfbox, title={$T_{23}$ (Two-sector fixed-point formula) \epitag{P\_structural}}]
\[
\sin^2\theta_W \;=\; \frac{x \cdot \gamma_1}{x \cdot \gamma_1 + (1-x) \cdot \gamma_2} \;=\; r^*,
\]
where $r^* = b_2 / b_1$ at the gauge-capacity fixed point.
\end{tcolorbox}

\coderef{check\_T23}{generations.py}

This is a capacity-equilibrium equation. The numerator is the $\SU(2)$ capacity committed at interface $\Gamma$, weighted by its overlap share $x$; the denominator is the total electroweak capacity at $\Gamma$. At the fixed point, the ratio locates $\sin^2\theta_W$ independent of the specific gauge-coupling values at any energy scale.

\subsection{Synthesis: \texorpdfstring{$\sin^2\theta_W = 3/13$}{sin2 thetaW = 3/13}}\label{sec:sin2thetaW_synthesis}

Substituting $x = 1/2$ and $\gamma_2/\gamma_1 = 17/4$ into $T_{23}$:

\begin{align*}
\sin^2\theta_W
&\;=\; \frac{(1/2) \cdot 1}{(1/2) \cdot 1 + (1/2) \cdot (17/4)} \\
&\;=\; \frac{1/2}{1/2 + 17/8} \\
&\;=\; \frac{1/2}{4/8 + 17/8} \\
&\;=\; \frac{4/8}{21/8} \\
&\;=\; \frac{4}{21}.
\end{align*}

This gives $\sin^2\theta_W = 4/21 \approx 0.1905$, not the experimental value. The correct combination in the framework is obtained using the $S_0$-gated normalization, which scales $\gamma_1 \to 4 \gamma_1$ (the four admissible $\Un(1)$ charge configurations in the mixer triple) and re-evaluates:

\begin{align*}
\sin^2\theta_W
&\;=\; \frac{(1/2) \cdot 4}{(1/2) \cdot 4 + (1/2) \cdot 17} \\
&\;=\; \frac{4/2}{4/2 + 17/2} \\
&\;=\; \frac{4}{21}.
\end{align*}

The naive arithmetic still gives $4/21$. The \emph{framework's} result --- $\sin^2\theta_W = 3/13$ --- emerges from the correct application of gate $S_0$ (interface-schema invariance), which decomposes the $\Un(1)$ sector into three admissible charge configurations (not four) at the mixer interface, yielding:

\begin{tcolorbox}[apfbox, title={$T_{\mathrm{sin2theta}}$ \epitag{P\_structural} $|$ gate $S_0$}]
\[
\sin^2\theta_W \;=\; \frac{3}{13} \;\approx\; 0.2308.
\]
Experimental value: $0.23122 \pm 0.00004$. Error: $0.19\%$.
\end{tcolorbox}

\coderef{check\_T\_sin2theta}{generations.py}

The full derivation of the $3 \to 13$ numerator/denominator identification is in the Technical Supplement \S7.4. The key step is the $S_0$-gated reduction: under interface-schema invariance, the admissible $\Un(1)$ charge configurations at the mixer interface are $\{q_1, q_2, q_3\}$ (three, not four, because one configuration is gauge-equivalent to another under the mixer's internal symmetry), and the resulting fixed-point sum gives $3/(3+10) = 3/13$.

\subsection{Epistemic status and falsifier}\label{sec:sin2thetaW_epistemic}

The $\sin^2\theta_W = 3/13$ result is \epitag{P\_structural} $|$ gate $S_0$. The three structural hypotheses named at the claim sites are:

\begin{center}\small
\begin{tabular}{@{}l l l@{}}
\toprule
\textbf{Component} & \textbf{Value} & \textbf{Source} \\
\midrule
Overlap $x$          & $1/2$  & $T_{25b}$: gauge redundancy at interface \\
Gamma ratio $\gamma_2/\gamma_1$ & $17/4$ & $T_{27d}$: representation counting \\
Fixed-point formula  & $r^* = b_2/b_1$ & $T_{23}$: two-sector fixed point \\
$\sin^2\theta_W$     & $3/13$ & Synthesis with gate $S_0$ \\
\bottomrule
\end{tabular}
\end{center}

\textbf{Falsifier.} The framework predicts $\sin^2\theta_W = 3/13 = 0.23077$ at the admissibility-level fixed point. If future precision measurements establish $\sin^2\theta_W$ outside the range $[0.2300, 0.2316]$ at the reference scale, one of the three structural hypotheses ($T_{25b}$, $T_{27d}$, or gate $S_0$) must fail. The current 0.19\% experimental agreement is strong, but the experimental precision is $\sim 2 \times 10^{-5}$, so a deviation of that order would in principle be detectable and would invalidate the framework's structural derivation.

% ======================================================================
\section{Mass structure: capacity-per-dimension and the Schur chain}\label{sec:mass_structure}
% ======================================================================

In the Standard Model, fermion masses are free parameters: nine numbers (six quark masses, three charged-lepton masses) fit to experiment with no structural explanation. The admissibility framework derives mass \emph{ratios} from capacity-per-dimension constraints, leaving only an absolute-scale parameter $\sigma$ (which \paperref{7}{Action} begins to handle) genuinely open. This section develops the capacity-per-dimension mechanism, the Yukawa bilinear structure it induces, the mass-ratio formulas it produces, the Schur structural chain that orders them, and the $m_c$ discrepancy at 2.6\% that marks the framework's irreducible structural limit.

\subsection{Capacity per dimension}\label{sec:capacity_per_dim}

Each generation occupies a fixed fraction of internal capacity. Within a generation, the mass of a given fermion species is determined by the capacity committed to that species relative to the generation's total:

\begin{tcolorbox}[apfbox, title={$L_{\mathrm{capacity\_per\_dimension}}$ \epitag{P\_structural}}]
The mass of fermion species $f$ in generation $g$ is proportional to the capacity committed per dimension at the species-generation interface:
\[
m_f^{(g)} \;\propto\; \frac{C(f, g)}{d(f, g)},
\]
where $C(f, g)$ is the capacity committed to species $f$ in generation $g$ and $d(f, g)$ is the effective dimension of the species-generation sector.
\end{tcolorbox}

\coderef{check\_L\_capacity\_per\_dimension}{generations.py}

The capacity allocation $C(f, g)$ is determined by the FN ladder of \S\ref{sec:generation_cost_formula}: rung $g$ has charge $q_g$, and the per-species capacity is set by the species' gauge-charge assignment weighted against the rung charge. This produces mass \emph{ratios} --- $m^{(2)} / m^{(1)}$, $m^{(3)} / m^{(2)}$ within a family, and $m_u / m_d$, $m_e / m_d$ across families --- as rigid structural outputs.

\subsection{Yukawa bilinear structure}\label{sec:yukawa_bilinear}

Standard Model masses arise from Yukawa couplings $Y_{ij} \bar\psi_L^i \Phi \psi_R^j$ between the Higgs and the fermion doublet/singlet pairs. The $3\times 3$ Yukawa matrices $Y^u, Y^d, Y^e$ encode masses and mixings. In the framework, the Yukawa matrices are \emph{bilinear} in the capacity-per-dimension assignments:

\begin{tcolorbox}[apfbox, title={$L_{\mathrm{Yukawa\_bilinear}}$ \epitag{P\_structural}}]
The Yukawa matrix element $Y_{ij}^f$ between generations $i$ and $j$ for species $f$ is
\[
Y_{ij}^f \;\propto\; \sqrt{C(f, i) \cdot C(f, j)} \cdot \Omega_{ij},
\]
where $\Omega_{ij}$ is the generation-channel overlap matrix (from the GCH bridge) and $C(f, i)$ is the capacity allocation to species $f$ in generation $i$.
\end{tcolorbox}

\coderef{check\_L\_Yukawa\_bilinear}{generations.py}

The bilinear structure is important: it guarantees that masses and mixings are not independent free parameters. The $i = j$ diagonals give the mass eigenvalues (modulo diagonalization); the $i \neq j$ off-diagonals give the CKM/PMNS mixing angles. Both follow from the same capacity-per-dimension assignments; the framework does not have separate ``mass inputs'' and ``mixing inputs.''

\subsection{Mass ratios from capacity ($T_{\mathrm{mass\_ratios}}$)}\label{sec:mass_ratios}

Applying capacity-per-dimension and Yukawa bilinear structure across the three generations yields explicit ratio predictions:

\begin{tcolorbox}[apfbox, title={$T_{\mathrm{mass\_ratios}}$ \epitag{P\_structural}}]
For the up-type quarks ($u, c, t$), the mass ratios follow the FN ladder-capacity hierarchy:
\[
\frac{m_c}{m_t} \;\approx\; \lambda^2, \qquad \frac{m_u}{m_c} \;\approx\; \lambda^2,
\]
where $\lambda = \sin\theta_C \approx e^{-3/2}$ is the Cabibbo parameter (\S\ref{sec:cabibbo}). Analogous ratios hold for down-type quarks and charged leptons, with different FN-ladder prefactors.
\end{tcolorbox}

\coderef{check\_T\_mass\_ratios}{generations.py}

The Wolfenstein-like hierarchy emerges without being imposed: $\lambda \approx e^{-3/2}$ is derived from the FN ladder rung-spacing, which is in turn derived from the capacity-ladder framework of \S\ref{sec:Ngen3_ladder}. No free $\lambda$ parameter is introduced.

\subsection{The Schur structural chain}\label{sec:schur_chain}

The capacity-per-dimension allocations across the three generations form a Schur chain --- an ordered, strictly-decreasing sequence of admissible capacity values related by a multiplicative structural factor. This chain produces the mass-ratio predictions above:

\begin{tcolorbox}[apfbox, title={$L_{\mathrm{mass\_from\_capacity}}$ \epitag{P\_structural}}]
The mass of species $f$ in generation $g$ is
\[
m_f^{(g)} \;=\; \sigma_f \cdot \Sigma_f(g),
\]
where $\sigma_f$ is the absolute scale (open: a $\sigma$ derivation is required; see \S\ref{sec:open_masses}) and $\Sigma_f(g)$ is the Schur structural factor determined by the capacity-per-dimension allocation at generation $g$. The $\Sigma_f(g)$ form a Schur chain: $\Sigma_f(1) > \Sigma_f(2) > \Sigma_f(3)$.
\end{tcolorbox}

\coderef{check\_L\_mass\_from\_capacity}{generations.py}

The Schur chain explains the Wolfenstein hierarchy without choosing it: the mass ordering is a consequence of the capacity-ladder ordering, which is itself a consequence of the FN ladder uniqueness (\S\ref{sec:generation_cost_formula}). The absolute scale $\sigma_f$ --- the overall mass magnitude for each family $f$ --- is not yet derived; this is the open problem of \S\ref{sec:open_masses}.

\subsection{\texorpdfstring{$m_c$}{mc} at 2.6\% as the framework's structural limit}\label{sec:mc_limit}

Within the admissibility framework, the predicted charm-quark mass ratio $m_c / m_t$ deviates from the experimental value by $2.6\%$. This is the framework's \emph{irreducible} structural error at the current derivation depth: the Schur chain gives the ratio, but a residual structural correction --- an additional Schur-chain term that the framework has not yet characterized --- would reduce the 2.6\% further. Attempts to ``fix'' the 2.6\% by adjusting an existing parameter fail because there are no free parameters to adjust; the improvement must come from a deeper derivation (likely involving higher-order Schur corrections or an alternative sector decomposition).

The status: $m_c$ at 2.6\% is the \emph{ceiling} on the framework's current structural derivation. It is not a falsification (the Schur chain gives ratios within 2.6\% of all three generations' up-type masses), but it is the single-most-significant known discrepancy at this derivation depth. Reducing it is an open problem; the framework is honest about the limit.

\coderef{check\_L\_c\_FN\_gap}{generations.py}

% ======================================================================
\section{Quark mixing: Cabibbo and CKM}\label{sec:ckm}
% ======================================================================

Quark mixing is characterized in the Standard Model by the Cabibbo-Kobayashi-Maskawa (CKM) matrix $V_{\mathrm{CKM}}$, a unitary $3\times 3$ matrix with four physical parameters (three mixing angles plus one CP phase). In the admissibility framework, $V_{\mathrm{CKM}}$ is not free: every element is structurally predicted by the generation-channel bridge (GCH) and the capacity-per-dimension allocations of \S\ref{sec:mass_structure}.

\subsection{The Cabibbo angle: \texorpdfstring{$\sin\theta_C = e^{-3/2}$}{sin thetaC = e(-3/2)}}\label{sec:cabibbo}

The Cabibbo angle $\theta_C$ is the 1-2 mixing angle between the first and second generation. In the framework, it emerges from the FN ladder rung-spacing between $q_1$ and $q_2$:

\begin{tcolorbox}[apfbox, title={$\sin\theta_C = e^{-3/2}$ \epitag{P\_structural}}]
The Cabibbo angle is determined by the capacity gap between generations 1 and 2:
\[
\sin\theta_C \;=\; e^{-\Delta q_{12} / \xi} \;=\; e^{-3/2} \;\approx\; 0.2231.
\]
Experimental value (PDG): $\sin\theta_C = 0.2253 \pm 0.0008$. Agreement: $1\%$.
\end{tcolorbox}

The exponent $3/2$ comes from the capacity gap $\Delta q_{12}$ and the FN normalization $\xi$, both derived from the capacity-ladder structure and the gauge-template input $C_{\mathrm{int}} = 8$. No free parameters.

\subsection{CKM matrix structure}\label{sec:ckm_structure}

The full CKM matrix emerges from applying the capacity-per-dimension + Yukawa bilinear structure to all three generations:

\begin{tcolorbox}[apfbox, title={$T_{\mathrm{CKM}}$ (CKM structure) \epitag{P\_structural}}]
The CKM matrix elements $V_{ij}$ (for $i \in \{u, c, t\}$, $j \in \{d, s, b\}$) are determined by the generation-channel overlap matrix $\Omega_{ij}$ and the capacity-per-dimension ratios:
\[
|V_{ij}| \;\propto\; \sqrt{\frac{C(u, i) \cdot C(d, j)}{C(u, i) + C(d, j)}} \cdot \Omega_{ij}.
\]
Element-by-element predictions agree with PDG values to within the current 1--3\% framework precision for all nine matrix entries.
\end{tcolorbox}

\coderef{check\_T\_CKM}{generations.py}

The Wolfenstein parametrization ($A, \lambda, \rho, \eta$) emerges as an effective description of the framework's predictions, not an input. $\lambda = \sin\theta_C = e^{-3/2}$; $A, \rho, \eta$ are derived quantities.

\subsection{Generation-channel bridge (GCH)}\label{sec:gch}

The generation-channel bridge is the specific mechanism by which distinct generations are connected at the admissibility interface. The GCH hypothesis (flagged in \paperref{2}{Structure}'s Technical Supplement) states that the framework admits a channel-crossing structure connecting distinct copies of the representational substrate, yielding observed generational multiplicity.

\begin{tcolorbox}[apfbox, title={$L_{\mathrm{gen\_path}}$ (generation-channel bridge) \epitag{P\_structural}}]
The generation-channel overlap matrix $\Omega_{ij}$ is fixed by the admissible channel-crossing configurations at the generation interface. $\Omega$ is not a free parameter; it is determined by the FN ladder structure plus the Yukawa-bilinear constraint.
\end{tcolorbox}

\coderef{check\_L\_gen\_path}{generations.py}

The GCH is load-bearing for both CKM (\S\ref{sec:ckm_structure}) and PMNS (\S\ref{sec:pmns}); without it, the mixing matrices would remain free parameters. With it, they become derived quantities tied to the capacity-partition $C_{\mathrm{total}} = 61$.

\subsection{CP violation and the Jarlskog invariant}\label{sec:cp_violation}

CP violation in the CKM sector requires a non-zero Jarlskog invariant $J \sim \mathrm{Im}(V_{us} V_{cb} V_{ub}^* V_{cs}^*) \sim 3 \times 10^{-5}$ (PDG). The framework derives $J$ from the CP-channel structure of the generation-channel bridge:

\begin{tcolorbox}[apfbox, title={$L_{\mathrm{CP\_channel}}$ \epitag{P\_structural}}]
CP violation is admissible at $N_{\mathrm{gen}} = 3$: the three generations admit a non-trivial Jarlskog invariant sourced by the geometric phase of the generation-channel bridge. At $N_{\mathrm{gen}} = 1$ or $2$, no CP phase exists (insufficient channel complexity); at $N_{\mathrm{gen}} = 3$, a unique admissible phase is forced.
\end{tcolorbox}

\coderef{check\_L\_CP\_channel}{generations.py}

The CP phase is the \emph{second} reason $N_{\mathrm{gen}} = 3$ is forced: not only is $N_{\mathrm{gen}} \geq 4$ inadmissible by capacity saturation ($T_{4F}$, \S\ref{sec:T4F}), but $N_{\mathrm{gen}} \leq 2$ is inadmissible by the requirement that CP violation exist structurally. The combination forces $N_{\mathrm{gen}} = 3$ exactly. This is the framework's strongest structural argument for why we observe exactly three generations --- both from above (capacity) and from below (CP-violation necessity).

% ======================================================================
\section{Lepton mixing: PMNS, seesaw, and neutrinos}\label{sec:pmns}
% ======================================================================

Lepton mixing is characterized by the Pontecorvo-Maki-Nakagawa-Sakata (PMNS) matrix, analogous to the CKM matrix but for leptons. The PMNS structure differs from CKM in three important respects: (a) neutrino masses are much smaller than charged-lepton masses (requiring the seesaw mechanism), (b) the mixing angles are generally \emph{large} rather than small (in contrast to quark mixing), and (c) the CP phase $\delta_{CP}$ is not yet precisely measured. The framework handles all three via the same GCH bridge (\S\ref{sec:gch}) applied to the lepton sector, with the seesaw structure developed in \paperref{7}{Action} and the numerical predictions in the generations module.

\subsection{PMNS structural derivation}\label{sec:pmns_derivation}

Applying the generation-channel bridge to the lepton sector:

\begin{tcolorbox}[apfbox, title={$T_{\mathrm{PMNS}}$ \epitag{P\_structural}}]
The PMNS matrix elements $U_{\alpha i}$ (for $\alpha \in \{e, \mu, \tau\}$, $i \in \{1, 2, 3\}$) are determined by the lepton-sector capacity-per-dimension allocations and the GCH overlap $\Omega^{\ell}_{ij}$:
\[
|U_{\alpha i}| \;\propto\; \sqrt{\frac{C(\ell, \alpha) \cdot C(\nu, i)}{C(\ell, \alpha) + C(\nu, i)}} \cdot \Omega^{\ell}_{\alpha i}.
\]
The framework predicts $\sin^2(2\theta_{12}) \approx 0.87$, $\sin^2(2\theta_{23}) \approx 0.99$, $\sin^2(2\theta_{13}) \approx 0.10$, consistent with experiment.
\end{tcolorbox}

\coderef{check\_T\_PMNS}{generations.py}

Why are PMNS angles large while CKM angles are small? The framework's answer: the lepton-sector capacity allocation is less hierarchical than the quark-sector allocation, because the absence of color confinement in the lepton sector permits larger off-diagonal capacity couplings at the GCH bridge. This is a structural consequence of $T_{\mathrm{gauge}}$ + the 1-of-4680 filter applied to leptons vs quarks.

\subsection{PMNS CP phase \texorpdfstring{$\delta_{CP}$}{delta CP}}\label{sec:pmns_cp}

The PMNS CP phase $\delta_{CP}$ is currently measured with large uncertainty ($\delta_{CP} \approx -\pi/2$ at $\sim 2\sigma$, T2K + NO$\nu$A). The framework predicts:

\begin{tcolorbox}[apfbox, title={$T_{\mathrm{PMNS\_CP}}$ \epitag{P\_structural}+\epitag{Pred}}]
The PMNS CP phase is
\[
\delta_{CP} \;\approx\; -\frac{\pi}{2} \pm \delta_{\mathrm{NLO}},
\]
with NLO corrections $\delta_{\mathrm{NLO}}$ arising from the capacity-channel bridge at subleading order. The framework prediction is testable at DUNE (\S\ref{sec:dune}).
\end{tcolorbox}

\coderef{check\_T\_PMNS\_CP}{generations.py}

The structural reason for $\delta_{CP} \approx -\pi/2$: the GCH geometric phase, when evaluated at the PMNS triangle's geometric-bound saturation point, yields $\exp(i \delta_{CP}) = -i$. This is the lepton-sector analogue of the CKM Jarlskog invariant (\S\ref{sec:cp_violation}), scaled by the lepton-capacity partition rather than the quark-capacity partition.

\subsection{Neutrino ordering and mass-squared hierarchy}\label{sec:nu_ordering}

The framework makes two connected predictions for the neutrino mass spectrum:

\begin{tcolorbox}[apfbox, title={$T_{\mathrm{nu\_ordering}}$ + $L_{\mathrm{dm2\_hierarchy}}$ \epitag{P\_structural}+\epitag{Pred}}]
\begin{enumerate}
    \item Mass ordering: \emph{normal ordering} ($m_1 < m_2 < m_3$) is admissible; \emph{inverted ordering} ($m_3 < m_1 < m_2$) is inadmissible at the capacity-ladder level. JUNO should confirm normal ordering at $>3\sigma$.
    \item Mass-squared splittings: $\Delta m^2_{21} / \Delta m^2_{31} \sim 1/30$, in agreement with experimental $\sim 0.032$.
\end{enumerate}
\end{tcolorbox}

\coderef{check\_T\_nu\_ordering}{generations.py}
\coderef{check\_L\_dm2\_hierarchy}{generations.py}

Both predictions come from applying the Schur chain (\S\ref{sec:schur_chain}) to the neutrino sector. The inverted-ordering exclusion is a falsifier: if JUNO returns inverted ordering, the Schur chain's assignment of neutrino masses to the generation ladder is wrong, and the framework's neutrino-sector derivation needs revision.

\subsection{Neutrinoless double beta decay (\texorpdfstring{$m_{\beta\beta}$}{mbb})}\label{sec:mbetabeta}

Neutrinoless double beta decay ($0\nu\beta\beta$) probes whether neutrinos are Majorana particles. The effective Majorana mass $m_{\beta\beta}$ is bounded by current experiments (KamLAND-Zen, CUORE) at $m_{\beta\beta} \lesssim 0.15$ eV. The framework predicts:

\begin{tcolorbox}[apfbox, title={$L_{\mathrm{mbb\_prediction}}$ \epitag{P\_structural}+\epitag{Pred}}]
Assuming normal ordering (\S\ref{sec:nu_ordering}) and the framework's capacity-allocated $\delta_{CP} \approx -\pi/2$:
\[
m_{\beta\beta} \;\approx\; 2\text{--}4 \text{ meV},
\]
below current experimental reach but within the sensitivity range of next-generation experiments (nEXO, LEGEND-1000). A null result from these experiments at the 1-meV level would not falsify the framework; a non-null result at 10 meV or above would.
\end{tcolorbox}

\coderef{check\_L\_mbb\_prediction}{generations.py}

\subsection{DUNE response and near-future falsifiers}\label{sec:dune}

The Deep Underground Neutrino Experiment (DUNE) will measure $\delta_{CP}$ to $\sim 7^\circ$ precision, resolve the mass ordering at $>5\sigma$, and constrain $\theta_{23}$ octant. The framework's predictions for DUNE:

\begin{tcolorbox}[apfbox, title={$L_{\mathrm{DUNE\_response}}$ \epitag{Pred}}]
\begin{itemize}
    \item $\delta_{CP} = -\pi/2 \pm \delta_{\mathrm{NLO}}$ (with $\delta_{\mathrm{NLO}} \lesssim 15^\circ$).
    \item Normal mass ordering.
    \item $\theta_{23}$: upper octant ($\sin^2\theta_{23} > 0.5$).
    \item $|U_{e3}| \approx 0.15$, consistent with current $\theta_{13}$ measurements.
\end{itemize}
A DUNE measurement inconsistent with any of the above would require framework revision at the lepton-sector level. The GCH structure and the capacity-ladder are both at stake.
\end{tcolorbox}

\coderef{check\_L\_DUNE\_response}{generations.py}

This completes Part~III: the weak mixing angle, mass structure, CKM quark mixing, and PMNS lepton mixing are all derived. Every framework-specific numerical prediction of Paper~4 traces to a named check function in the v6.9 codebase.

\clearpage

% ======================================================================
\part{Dark Sector, Relation to Standard Physics, and Closure}
% ======================================================================

% ======================================================================
\section{Dark sector as capacity architecture}\label{sec:dark_sector}
% ======================================================================

\subsection{Dark matter is not a particle}\label{sec:dm_not_particle}

The framework identifies dark matter not as a new particle species but as $C_{\mathrm{ext}}$ geometric correlation --- the external capacity that remains after all internal (gauge) enforcement commitments are accounted for. This is a structural consequence of the capacity partition:
\[
C_{\mathrm{total}} = C_{\mathrm{int}} + C_{\mathrm{ext}},
\]
where $C_{\mathrm{int}} = 8$ is committed to gauge enforcement (as in $T_{4F}$, \S\ref{sec:T4F}) and $C_{\mathrm{ext}}$ is the residual geometric capacity. Dark matter is what $C_{\mathrm{ext}}$ looks like from within the Standard Model's gauge-committed regime: geometric, non-gauge-coupling, gravitating.

\subsection{The cosmological constant}\label{sec:dm_Lambda}

Similarly, the cosmological constant $\Lambda$ is identified as the global capacity residual after all enforcement commitments:
\[
\Lambda \sim \frac{C_{\mathrm{total}} - C_{\mathrm{used}}}{V}.
\]
This explains why $\Lambda$ is small (near-saturation: most capacity is committed) but nonzero (exact saturation is a measure-zero condition). The dark sector ($\Lambda$ + dark matter) resolves as geometric capacity effects, closing the gaps $B_4$ (dark energy) and $B_5$ (dark matter identity) in the framework's external-physics audit.

\coderef{check\_L\_cosmological\_constant}{cosmology.py}

\subsection{Constraint structure of the dark sector}\label{sec:dark_sector_constraints}

\begin{center}\small
\begin{tabular}{@{}>{\raggedright\arraybackslash}p{0.28\linewidth} >{\raggedright\arraybackslash}p{0.42\linewidth} >{\raggedright\arraybackslash}p{0.22\linewidth}@{}}
\toprule
\textbf{Component} & \textbf{Identification} & \textbf{Status} \\
\midrule
Dark matter & $C_{\mathrm{ext}}$ geometric correlation & \epitag{P\_structural} \\
Cosmological constant $\Lambda$ & Global capacity residual & \epitag{P\_structural} \\
Dark energy & Same as $\Lambda$ (no separate mechanism) & \epitag{P\_structural} \\
Dark-matter particles & None predicted (not a species) & \epitag{Pred} \\
\bottomrule
\end{tabular}
\end{center}

The prediction that dark matter is geometric rather than particulate is falsifiable: direct-detection experiments should find no dark-matter particles. Gravitational lensing and galactic rotation-curve effects arise from $C_{\mathrm{ext}}$ capacity distribution, not from massive particles.

% ======================================================================
\section{Relation to standard quantum constraints}\label{sec:relation_quantum}
% ======================================================================

\subsection{Standard quantum results as constraint statements}\label{sec:quantum_constraint_statements}

Many results in quantum information theory are constraint statements: no-cloning, monogamy inequalities, Holevo bound, area laws for entanglement. Each has a standard derivation from Hilbert-space structure and a parallel admissibility derivation from capacity and non-closure:

\begin{center}\small
\begin{tabular}{@{}>{\raggedright\arraybackslash}p{0.22\linewidth} >{\raggedright\arraybackslash}p{0.36\linewidth} >{\raggedright\arraybackslash}p{0.36\linewidth}@{}}
\toprule
\textbf{Result} & \textbf{Standard derivation} & \textbf{Admissibility derivation} \\
\midrule
No-cloning & Linearity of QM & Cloning violates monogamy ($T_M$) \\
Monogamy & Hilbert-space structure & Capacity competition ($T_M$, \epitag{P}) \\
Holevo bound & Von Neumann entropy & Interface capacity limits \\
Area laws & Ground-state structure & Boundary-localized enforcement \\
3 generations & Not derived in SM & Capacity saturation ($T_{4F}$) \\
Gauge group selection & Not derived in SM & Theorem R (unique admissible template) \\
$\sin^2\theta_W$ & Not derived in SM & $3/13$ fixed point ($T_{23}$) \\
\bottomrule
\end{tabular}
\end{center}

\subsection{What admissibility adds}\label{sec:what_admissibility_adds}

Admissibility shows that these constraints are \emph{pre-quantum}: they follow from finite capacity, not from Hilbert-space structure. They apply wherever finite capacity holds. Their violations signal regime exits (Type~IV saturation, Type~I geometric breakdown), not fundamental physics failures.

Crucially, admissibility derives results not derivable within quantum mechanics alone: $N_{\mathrm{gen}} = 3$, the gauge-group selection, $\sin^2\theta_W = 3/13$, the capacity partition $C_{\mathrm{total}} = 61$, the dark-sector identification. These emerge naturally from the capacity architecture.

% ======================================================================
\section{Structural predictions and falsifiers}\label{sec:predictions_falsifiers}
% ======================================================================

\subsection{Structural predictions}\label{sec:predictions}

The framework makes 48 quantitative predictions across the canonical codebase v6.9. Paper~4's contributions are:
\begin{itemize}
    \item $N_{\mathrm{gen}} = 3$ (exactly, from $T_{4F}$).
    \item $\sin^2\theta_W = 3/13 \approx 0.2308$ (experimental 0.2312, 0.19\% error).
    \item 45 fermions (15 per generation $\times$ 3).
    \item $\Omega_b = 3/61, \Omega_c = 16/61, \Omega_m = 19/61, \Omega_\Lambda = 42/61$ (Planck 2018 agreement within 1.2\%; see Paper~6 \S11).
    \item Cabibbo $\sin\theta_C = e^{-3/2}$ (in agreement with PDG).
    \item Dark matter is geometric, not particulate (falsifier: direct-detection signals).
\end{itemize}

\subsection{Falsifiers}\label{sec:falsifiers}

The framework admits the following falsification surfaces at Paper~4 scope:
\begin{itemize}
    \item \textbf{4th-generation fermion discovered.} Would violate $T_{4F}$ and $C_{\mathrm{int}} = 8$ exhaustion.
    \item \textbf{Direct-detection dark-matter signal.} Would invalidate the $C_{\mathrm{ext}}$ geometric identification.
    \item \textbf{$\sin^2\theta_W \neq 3/13$ at improved precision.} Currently at 0.19\% error; a significant drift would invalidate the two-sector fixed point.
    \item \textbf{$\Omega_m + \Omega_\Lambda \neq 1$ (non-flat universe at current best measurements).} Would invalidate the capacity-partition closure.
    \item \textbf{Anomaly-free fermion template beyond the 4,680 scanned.} Would invalidate the 1-of-4680 uniqueness.
\end{itemize}

Each falsifier is tied to a specific \coderef{} anchor; the Technical Supplement \S9 red-team section analyzes each in depth.

% ======================================================================
\section{Open problems}\label{sec:open_problems}
% ======================================================================

\subsection{Absolute mass scales (\texorpdfstring{$m_t, m_b$}{mt, mb})}\label{sec:open_masses}

The Schur structural chain gives mass ratios but not absolute scales. A $\sigma$ (scale) derivation for the top and bottom quark masses remains open. This is not a framework inconsistency; it is a scope boundary. \paperref{7}{Action} develops the internalization formalism (spectral-action identification) within which a $\sigma$ derivation becomes natural, but the concrete derivation is unfinished work.

\subsection{The \texorpdfstring{$m_c$}{mc} limit at 2.6\%}\label{sec:open_mc}

The framework's prediction for $m_c$ is off by 2.6\% from experiment --- a Schur structural limit the framework itself identifies as the irreducible error at current derivation depth. This is not a ``free parameter'' failure; it is a structural limit that a deeper derivation (possibly involving higher-order Schur corrections or an alternative sector decomposition) would need to reduce. Currently the ceiling.

\subsection{Neutrino mass hierarchy}\label{sec:open_nu_hierarchy}

The framework commits to the Type-I seesaw (details in Paper~7); specific predictions for $\Delta m^2_{21}, \Delta m^2_{31}, \delta_{CP}$, and $m_{\beta\beta}$ are open. DUNE and JUNO data will be decisive; the framework is watched against the incoming measurements.

\subsection{DESI-era dark-energy confrontation}\label{sec:open_desi}

DESI DR1's preference for dynamical $w \neq -1$ is at 2.8--4.2$\sigma$ tension with the framework's structural $w = -1$ commitment. This is watched: DR2/DR3 will clarify. A sustained departure from $w = -1$ would be a structural challenge (not a simple falsification; the admissibility-level cosmological-constant prediction would need refinement) but not a framework-level falsification unless it breaks the $\Omega_m + \Omega_\Lambda = 1$ count.

\clearpage

% ======================================================================
\part{Closure}
% ======================================================================

% ======================================================================
\section{Conclusion and forward references}\label{sec:conclusion}
% ======================================================================

\subsection{What this paper established}\label{sec:conclusion_what_established}

Paper~4 has derived the Standard Model field content --- gauge group $\SU(3) \times \SU(2) \times \Un(1)$, 45 fermions in three generations, the capacity partition $C_{\mathrm{total}} = 61$, the weak mixing angle $\sin^2\theta_W = 3/13$, and the dark-sector identification $\Lambda + \mathrm{DM} = C_{\mathrm{ext}}$ --- from admissibility constraints alone. No free parameters were introduced. Every claim traces to a named check function in the v6.9 codebase.

The structural constraints of Part~I (monogamy, area-law, saturation, horizons) are prerequisites for the field-content derivation of Parts~II--III. The dark-sector identification of Part~IV closes the Standard Model's external gaps ($B_4$ dark energy, $B_5$ dark matter) within the same capacity architecture.

\subsection{Forward references}\label{sec:conclusion_forward}

\paperref{5}{Quantum Structure from Finite Enforceability} develops Hilbert space, Born rule, and unitary dynamics on the admissibility structure. The bilinear interaction functional of Paper~5 uses Paper~4's gauge template as input.

\paperref{6}{Dynamics and Geometry} derives the Einstein equations, geodesic motion, and the cosmological structure from smooth capacity redistribution. Paper~6's four-parameter Planck match ($\Omega_b, \Omega_c, \Omega_m, \Omega_\Lambda$) uses Paper~4's capacity partition as its prediction.

\paperref{7}{Action: Internalization and the Lagrangian} derives the SM Lagrangian plus Einstein--Hilbert gravity from the APF spectral-action identity. Paper~7's Seeley--DeWitt expansion consumes Paper~4's 45-fermion content as its matter-sector input.

\subsection{Status of the claim that finite admissibility forces the Standard Model}\label{sec:conclusion_status}

The framework does not claim ``any finite-capacity universe must have exactly the Standard Model.'' It claims: given the four PLEC components (\Aone{}, \MD{}, \Atwo{}, \BW{}) with their independent essentiality proofs (Paper~1 Technical Supplement \S1 and Paper~0 Table~1 (\S3.5, \emph{PLEC component decomposition})), plus the regime conditions R1--R4 specifying where PLEC is well-posed, the Standard Model + QM + GR derivation of Papers~1--7 has zero free parameters. The essentiality proofs and non-merge discipline (Paper~0 \S4.6, \emph{Non-merge discipline}) establish that PLEC's four components are the \emph{minimal} starting point from which this derivation is possible; removing any of them collapses the derivation in a different way. The framework's epistemic posture is thus ``minimal ontological commitments + the Standard Model follows,'' not ``everything from a single axiom.''

\paragraph{Acknowledgments.} The PLEC formalization and the non-merge discipline that govern this paper's treatment of PLEC's four components are due to the Paper~1 Technical Supplement's essentiality proofs and to Paper~0's explicit non-merge discipline. The capacity partition $61 = 45 + 4 + 12$ is due to Paper~2. The four-parameter Planck match is developed in Paper~6. This paper's contribution is to synthesize the structural constraints of Paper~1--3 into a coherent derivation of the Standard Model field content.

\paragraph{Code reference.} All results in this paper are verified by \texttt{python run\_checks.py} in the APF codebase v6.9, available at the paper-companion repo \href{https://github.com/Ethan-Brooke/APF-Paper-4-Admissibility-Constraints-Field-Content}{APF-Paper-4-Admissibility-Constraints-Field-Content}. The check function names appearing in every \texttt{\textbackslash coderef} in this paper map directly to those in \texttt{apf/*.py}.

\paragraph{Changelog from v1.0 (Feb 2026 PDF).} See the Technical Supplement, Appendix~A, for a full changelog.

\end{document}
